In this chapter, I turn to the second facet of multifaceted geospatial data: its \emph{ensemble variability}.
%
Climate projections rarely arrive as a single forecast; instead, large ensembles of model runs each trace a plausible spatiotemporal future, varying both from the systematic forcings scientists wish to study and from the chaotic internal variability of the climate system.
%
Conventional statistical summaries such as spatial means and variances compress this richness, often obscuring the very spatial and temporal structures that distinguish one plausible future from another.
%
To reconcile this, I present \emph{ClimateSOM}, a visual analysis workflow that abstracts the members of a spatiotemporal ensemble into a distribution over a self-organizing-map-defined, user-steerable 2D space, supporting the exploration, comparison, and clustering of model behavior with LLM-assisted sensemaking.
%
I evaluate ClimateSOM through case studies on LOCA-downscaled CMIP6 projections over the western United States and an expert review with domain scientists.
%
The workflow lets an analyst explore and \emph{compare} the behavior of model runs, but it leaves a central question unanswered: whether the comparisons it surfaces are statistically significant or merely apparent.
%
To this end, I extend the workflow from exploration to inference.

In the second section of this chapter, I extend ClimateSOM with a statistical testing schema built on the Earth Mover's Distance, aiming to determine---with statistical significance---whether the \emph{distribution} of spatial precipitation patterns, and the co-variability between regions, is expected to change under a changing climate.
%
I demonstrate this extension on the co-variability between California's Sacramento Bay-Delta Watershed and U.S. West Coast precipitation, where I find a strengthening departure from historical relationships that I term ``decoherence.''
%
The remainder of this chapter develops the motivation, design, and assessment of both the workflow and its inferential extension in detail.

\section{ClimateSOM --- The workflow}\label{sec:climatesom-workflow}
\subsection{Introduction}\label{sec:intro}
% \subsubsection*{\note{Proposed Structure}}
% \begin{tightEnumerate}
%     \item As climate models increase in complexity (e.g. CMIP6 is significantly more complex than CMIP5) there is an increasing urgency to understand their outputs.
%     \item Recent climate ensemble data like CMIP6 contains within it many Global Climate Models (GCM) and Shared Socioeconomic Pathways (SSP), each pair arising in a similar yet dissimilar spatiotemporal time series over some given spatial domain.
%     \item Understanding each member's output, comparing different member's outputs, and clustering these results is increasingly central to climate science and understanding future variability of the climate.
%     \item These three tasks are the main concern - explore one, compare two, cluster many
%     \item However the current state-of-the-art analytic methods founded in climate science to address these tasks are statistic-based methods (e.g. computing and comparing spatial means over different SSPs) and while they have proven to be successful \fix{cite} leave room for improvement.
%     \item Similarly, state-of-the-art tools in visual analytics falls short in its interpretability for domain scientists.
%     \item To this end, we propose a visual analysis workflow to equip climate scientist with a visual tool and language to address these questions.
%     \item With \ClimateSOM we apply a self-organizing map (SOM), a well recognized analysis tool in climate science, with carefully tuned parameters, along with dimensionality reduction and annotation, to abstract spatiotemporal time series to a distribution over a user-steerable space. 
%     \item Finally we demonstrate that \ClimateSOM can produce scientific insights about climate ensemble data that are otherwise time consuming and difficult to yield otherwise.
% \end{tightEnumerate}

%% As climate models continue to increase in complexity, the need to understand their outputs has become more urgent. 
%As climate models and their ensemble data sets become increasingly complex, the need to understand their outputs has grown. \fix{cite climate science papers}
%%
%One such example is the transition from Coupled Model Intercomparison Project (CMIP) CMIP5 to CMIP6 \cite{eyring2016overview}, which represents a significant leap in the size of the ensemble.
%%
%The CMIP6 ensemble dataset contains a wide array of outputs, including data from numerous Global Climate Models (GCMs) and Shared Socioeconomic Pathways (SSPs). 
%%
%These outputs form spatiotemporal time series that, while arising from similar processes, can crucially vary significantly across different runs.

%The complexity of this data necessitates a deeper understanding of the individual climate model outputs as well as the relationships between them. 
%%
%% In particular, we identify three key motivating tasks for ensemble datasets - explore one, compare two, and cluster many.
%In particular, as previously presented by Wang in their survey of ensemble data visualization \cite{wang_visualization_2019}, we also identified three key tasks for working with ensemble datasets: exploring individual models, comparing results across models, and clustering multiple outputs.
%%
%% As a result exploring the outputs of each GCM-SSP pair, comparing results from different models, and clustering these results have become critical components of climate science \cite{almazroui2020projections}. \fix{add other CMIP6 paper}
%%
%This is particularly important as researchers aim to better understand the variability of future climate scenarios.

%However, current state-of-the-art analytical methods in climate science, which often rely on statistical techniques such as computing and comparing spatial means across different SSPs, while successful \fix{add some papers to this end}, leave room for improvement in terms of providing deeper insights.
%%
%The field of visual analytics has also proposed many methods to this end \cite{kappe_exploring_2019}, \fix{cite more}; however, they short in its ability to support the range of tasks—such as exploration, comparison, and clustering of ensemble models—that \ClimateSOM is designed to address.

%%
%To address these challenges, we propose a novel visual analysis workflow designed to equip climate scientists with an intuitive tool and language for exploring and analyzing climate ensemble data.
%%
%In particular, we introduce \ClimateSOM, a workflow that applies self-organizing maps (SOMs)—a well-established analytical tool in climate science—in combination with carefully tuned parameters, steerable dimensionality reduction, and annotation. 
%%
%In all, this methodology abstracts spatiotemporal time series into a distribution over a user-steerable space, enabling more flexible and insightful exploration of the data.

%Finally, we demonstrate that \ClimateSOM can yield scientific insights from climate ensemble data that would otherwise be time-consuming and difficult to extract, thus offering a valuable advancement in the field of climate science.
% \begin{tightItemize}
%     \item \note{Spatiotemporal ensemble data is increasingly large and requires solutions to explore, compare, and cluster them}
%     \item \note{Despite work in this area, its still unclear how we can communicate the dynamics of spatiotemporal data in a concise way while also supporting explicit support for other tasks like comparison and clustering.}
% \end{tightItemize}

Ensemble datasets are increasingly integral to scientific research as computing resources have increased in power \cite{wang_visualization_2019, kehrer_visualization_2013}. 
Many disciplines, from computational fluid dynamics~\cite{patchett2017deep} to climate science~\cite{maher2021large}, rely on ensemble datasets to capture and predict variability by running their models under different conditions and parameters.
The result is a collection of model runs---an \textit{ensemble} dataset---that supports exploration of possible outcomes, sensitivity analysis, and uncertainty quantification. However, working with ensemble datasets remains difficult due to their sheer size and complexity~\cite{wang_visualization_2019}.
% The output of such a process, a set of many model runs known as an \textit{ensemble} dataset, helps researchers explore the range of possible outcomes, understand the sensitivity of the outputs with respect to the model variants or parameters, and quantify uncertainty in predictions. Despite its ubiquity, a major challenge with working with ensemble datasets is often their sheer size and complexity \cite{wang_visualization_2019}.

% Since a single model run often requires a specialized analysis, this difficulty only grows when dealing with an entire ensemble.
% Since analyzing a single output set from a particular  model run often requires individual analysis, dealing with an entire ensemble only amplifies this difficulty.
% They often exhibit substantial variability both within and across model runs, making them difficult for domain scientists to fully interpret them.
% ensemble datasets often highly vary within and across model runs, 

Climate science in particular relies on ensemble datasets to examine plausible climate futures under varying conditions, such as greenhouse gas and aerosol emissions or land use changes
~\cite{murphy_quantification_2004, eyring2016overview, maher_max_2019}.
%
Each model run often represents  a spatiotemporal time series, where each time step is a climate outcome (e.g. precipitation) over some spatial regions.
%
These model runs may come from different global climate models~\cite{eyring2016overview, maher_max_2019} and use varying input assumptions, producing a broad set of outcomes. Understanding and communicating the resulting variability and patterns remains a major challenge.


Although existing analyses in climate science provide valuable insights---often through statistical summaries like spatial means and variances across model runs---they may overlook key spatial and temporal structures~\cite{kappe_exploring_2019}.
Visual analytics techniques have been developed to support interpretation of ensemble data~\cite{biswas2016visualization, kappe_exploring_2019, evers_interactive_2023}, but many struggle to provide both interpretability and robustness when dealing with highly complex, chaotic model behavior.
%
% The main challenge, particularly with climate ensemble datasets, as previously identified\cite{potter2009visualization}, is that model runs are often both highly chaotic and complex.
%
% In addition to the systematic variations introduced by the model inputs that climate scientists are interested in exploring, they usually also include unknown amounts of random variation.
% In addition to the systematic climate variations driven by the model inputs that are sought after, any analytics and must contend with unknown amounts of random variation.
% they include unknown amounts of random variation in addition to the systematic variations of interest introduced by the GCMs and SSPs.
%
% In face of such variability and complexity, 
%
% As Kappe outlines, in the context of ensemble spatiotemporal simulations, ``direct visualizations \textit{(of climate ensembles)} are prohibitive with respect to space, time, and mental capacity of the user'' \cite{kappe2022visual}. 
%
% Simple exploratory visualizations of the this data often reveal that the innumerable number of variations observed in the dataset is too complex to understand or visualize in its entirety.
In this context, even summarizing a single model run clearly is difficult, and comparing or clustering multiple runs poses a greater challenge.


% As a result, most approaches resort to \textit{Summary statistics} 

To overcome these limitations, we introduce a novel visual analysis workflow that equips climate scientists with an intuitive tool and visual language for exploring and analyzing ensemble climate projection datasets. 
%
In line with findings by Wang et al. in their comprehensive survey of ensemble data visualization \cite{wang_visualization_2019}, we identify three essential tasks when working with such datasets: \textit{(1) exploring a single set of model run(s)}, \textit{(2) comparing the behavior across two sets of model runs}, and \textit{(3) clustering the behavior of multiple model runs}. 
%
These tasks are crucial for uncovering insights into climate variability and making meaningful inferences about these projections. 
%
These tasks along with its interpretability serves as a guiding principle for \ClimateSOM.

The overarching goal for \ClimateSOM is data abstraction by transforming spatiotemporal time series into a distribution over a user-steerable 2D space.
%
Central to \ClimateSOM is a self-organizing map (SOM)\cite{kohonen_self-organizing_1990}, a widely recognized analytical tool in climate science, along with steerable dimensionality reduction and annotations.
%
% Expanding on past work on SOM training, we highlight the impact of SOM training parameters and its importance to the success of this visual workflow.
%
Through this transformation, \ClimateSOM enables flexible, interpretable exploration of diverse model behavior, helping researchers detect patterns and relationships that would otherwise be difficult to uncover.

We further integrate LLMs to assist in sensemaking of the SOM-defined 2D space.
%
\ClimateSOM introduces an LLM-empowered bidirectional query methods that allows users to (1) locate regions of interest in the 2D space based on natural language queries and (2) generate textual summaries of selected regions.
%
We demonstrate the utility of \ClimateSOM through case studies using LOCA-downscaled CMIP6 climate projections~\cite{pierce2023future} across two regions in the western United States. In addition, we conduct an expert review with six domain scientists to assess the system’s usability and its ability to support scientific insight.
%
In short, the contributions of this work are the following:
\begin{itemize}[noitemsep, topsep=0pt]
    \item A novel workflow for visual analysis of a spatiotemporal ensemble that abstracts the ensemble members as a distribution over a steerable 2D space framed by a self-organizing map.
    \item Integration of LLMs to assist in sensemaking of the SOM-defined 2D space in which \ClimateSOM operates.
    \item \ClimateSOM, a visualization interface that supports this workflow end-to-end, enabling users to explore, compare, and cluster ensemble model runs.
    \item Case studies on the LOCA-Downscaled CMIP6 climate model projections over two different regions of the Western United States and interviews with six domain experts to evaluate its efficacy.
        % to demonstrate the utility of \ClimateSOM in uncovering insights from climate ensemble datasets.
\end{itemize}



\subsection{Related Works}\label{sec:relworks}
%In their survey, Wang et al. identifiy five \textit{dimensions} of ensemble data: variable, location, time, member, and ensemble. 
%%
%They also define six \textit{analytic tasks} for ensemble data: overview, compare, cluster, trend, feature, and parameters.
%%
%Each \textit{analytic task}, when applied to different \textit{dimensions}, demands a distinct analysis and visualization strategy\cite{wang_visualization_2019}.
% For any of the \textit{analytic tasks}, the \textit{dimension} considered brings rise to a entirely different analysis and visualization strategy \cite{wang_visualization_2019}.
%
% We adopt this taxonomy to structure our discussion of the visualization strategies for ensemble climate data.

\subsubsection{Visual Analytics for Climate Ensemble Data}
% \subsubsection*{\note{Proposed Structure}}

% \begin{tightItemize}
%     \item Here are past work that have tried climate ensemble data visualization.
%     \item Given the complexity of this data, visualization and analytics of such data often rely on data abstraction and aggregation to compress the data into a concise form.
%     \item For climate ensemble data, the most straightforward intermediate data abstraction are summary statistics as is populate in climate science for its simplicity and interpretablilty\cite{kehrer_brushing_2010,hollt_ovis_2014}. However, past work have also considered graph-based approach, where a graph is constructed to represent the ensemble \cite{shu2016ensemblegraph, kappe_analysis_2019}, topology-based methods, where isocontours, spaghetti plots or similar methods are employed to represent the ensemble data while retainig the spatial context,  \cite{kappe_topology-based_2019,kappe_topology-based_2022, liu_visualizing_2015, vietinghoff_visual_2021, sanyal_noodles_2010, ma_interactive_2019}, or even as 2D images\cite{ise_varenn_2020}.

% \end{tightItemize}

% Given the abundance of ensemble data sets from large-scale simulations, many visualization techniques to investigate the ensemble data have been proposed in literature \cite{potter2009ensemble}.

The increasing availability of ensemble data from large-scale climate simulations has led to the development of a range of visual analytics methods designed to manage and interpret uncertainty, spatial-temporal patterns, and inter-model variability~\cite{kehrer_brushing_2010, biswas_visualization_2017, ferstl_time-hierarchical_2017, wang_multi-resolution_2017,biswas2016visualization, kappe_exploring_2019}. 
%
% Given the complexity of this data, visualization and analytics of such data often rely on data abstraction and aggregation to compress the data into a concise form.
%
% For climate ensemble data, the most common data abstraction is summary statistics as is popular in climate science for its simplicity and interpretablilty \cite{kehrer_brushing_2010,hollt_ovis_2014}. 
%
% Of the many methods, statistics~\cite{kehrer_brushing_2010, hollt_ovis_2014} and matrix decomposition~\cite{chandler_characterizing_2024} have been central tools for some previous works.
%
Past works have considered topology-based methods, where isocontours, spaghetti plots or similar methods are employed to preserve spatial context while encoding ensemble variability~\cite{kappe_topology-based_2019,kappe_topology-based_2022, liu_visualizing_2015, sanyal_noodles_2010, ma_interactive_2019}.
Others have proposed graph-based representations of ensemble relationships~\cite{shu2016ensemblegraph, kappe_analysis_2019} or treated ensemble members as collections of images to facilitate comparative analysis~\cite{ise_varenn_2020}.
%
% In \cite{shu2016ensemblegraph} Shu et al. presented EnsembleGraph, a visual analytics tool for time-varying ensemble data that abstracted spatial regions as nodes to develop a graph-based representation of the ensemble.
%
Kappe et al. introduced a Sankey-diagram-based visualization using \textit{k}-means clustering on the entire climate field to summarize variability in decadal predictions~\cite{kappe_exploring_2019}, while Evers et al. employed spatial correlation measures and multidimensional scaling (MDS) to embed ensemble members for comparative visualization~\cite{evers_uncertaintyaware_2021}.
%
% This work was extended replacing hierarchical clustering with \textit{k}-means clustering \cite{kappe_analysis_2019}.
% to build a more advanced visual analytics tool 
%
% More recently, Evers et al. demonstrated a visualization workflow for climate ensemble via calculating spatial correlation across the spatial domain and then embedding the climate ensemble using multi-dimensional scaling (MDS)~\cite{evers_uncertaintyaware_2021}.
%
% This method allows users to understand regions of spatial domain that exhibit similar behavior across the ensemble members.
%
% This work was further extended for multi-field climate ensembles by calculating spatial correlations on a concatenated time series across all fields~\cite{evers_interactive_2023}.

%In both set of works, clustering of the ensemble data is identified as a key \textit{analytic task}, though they differ the \textit{dimension}; the former clusters the ensemble members, while the latter clusters the spatial domain.
%%
%\ClimateSOM aligns closer to \cite{kappe_analysis_2019, kappe_exploring_2019} where the clustering on the ensemble members is the primary focus.

% Given the complexity of climate ensemble data, the balance between the level of abstraction and the amount of information retained is a key challenge in its visualization, and we propose \ClimateSOM as another approach to address this challenge, this time to abstract the data into a 2D distribution.

However, in the context of visualizing the overall behavior of a single climate ensemble model run or its comparison, past works have mostly relied on summary statistics or what Kappe refers to the as the \textit{mean field approach}~\cite{kappe2022visual}.
%
As Kappe points out, however, this aggregation can be misleading since climate ensemble data seldom vary around a single mean, instead exhibiting multiple means of varying importance and distributional characteristics.

%
Departing from this \textit{mean field} aggregation is a key motivation for \ClimateSOM. Rather than reducing the ensemble to a single statistical summary, our approach models it as a 2D spatially organized distribution, preserving the variability across model runs. We use a self-organizing map (SOM) to maintain local topological relationships while revealing fine-grained differences between ensemble members. In contrast to prior systems, which typically support only isolated subtasks, such as single-run exploration or pairwise comparison, \ClimateSOM provides integrated support for exploration, comparison, and clustering. This unified capability addresses a gap in existing climate ensemble visual analytics.

%\subsection{Dimensionality Reduction (DR) for Spatiotemporal Data}
%The manifold hypothesis \cite{fefferman2016testing} that posits that high-dimensional data lies on a low-dimensional manifold has been a driving force behind the development of dimensionality reduction (DR) techniques.
%%
%Notato and Aupetit provide a comprehensive survey of DR or projection methods for visual analytic tasks \cite{nonato_multidimensional_2019}.
%%
%Spatial or spatiotempoiral data is no exception to this and DR methods can be readily to spatial data as well, providing an avenue to visually explore such data.
%%
%Classical Multidimensional Scaling (MDS) is often used in these works \cite{piccolotto_undrground_2024} \note{Cite more ST-MDS works.}
%%
%In \cite{piccolotto_undrground_2024}, Piccolotto et al. used MDS to visualize latent components and their relationships that were extratcted from a Spatial Blind Source Separation (SBSS) analysis.

%%
%As Nonato and Aupetit outline the in any DR or similar projection method, the choice of the dissilarity measure is crucial to the success of the method.
%%
%Although Euclidean distance is the most common choice, its expressive power is limited in cases where the distance is large as in images \cite{agrawal2021minimum}.
%%
%Kappe explored the role of dissimilarity measures in the context of climate data \cite{kappe2022visual}

%Brehmer et al. identified a number of common tasks for visualization with dimensionality reduction, including the intepretation of the synthesized dimension \cite{brehmer_visualizing_2014}.



\subsubsection{SOMs in Visual Analytics and Climate Science}
Self-organizing maps (SOMs)~\cite{kohonen_self-organizing_1990} are commonly used for their ability to project high-dimensional data onto a structured, low-dimensional grid while preserving local topological relationships. This pseudo-continuous 2D layout enables spatial interpretation and supports visual pattern recognition.
In visual analytics, SOMs have been used to structure time-series and spatial data, as demonstrated by Andrienko et al. in their analysis of crime patterns ~\cite{andrienko_space-time_2010}. 
Beyond this, SOMs have been applied in diverse domains such as astrophysics and cytometry, where the preservation of local structure facilitates cluster identification and improves interpretability~\cite{chen_self-organizing_2013, hemmati_bringing_2019, sacha_somflow_2018, chushig-muzo_data-driven_2020, van_gassen_flowsom_2015,quintelier_analyzing_2021, tu_phrasemap_2024, eirich_irvine_2022}.

In climate science, SOMs have often been used for pattern extraction, where extracted nodes represent particular spatial patterns that may be linked to known physical phenomena \cite{liu_performance_2006}.
% The ``\textit{self-organizing}'' property of SOMs, which enables dimensionality reduction, is often credited as a key feature that makes analysis via SOMs fruitful in Climate Sciences, where extracted nodes and patterns can often be associated with known physical phenomena \cite{liu_performance_2006}.
%
For instance, Reusch et al. used SOMs to identify the North Atlantic Oscillation (NAO) in climate data \cite{reusch_north_2007}, while Mihanovi\'{c} et al. used SOMs to identify surface circulation patterns in the northern Adriatic Sea \cite{mihanovic_surface_2011}.
%
% Nourani et al. used SOMs to preprocess precipitation satellite data for further analysis \cite{nourani_using_2013} and Loikith et al. used SOMs to identify large-scale meteorological patterns (LSMPs) over the Northwestern United States \cite{loikith_characterizing_2017}.
%
For more applications of SOMs in climate science or related areas, we refer the reader to \cite{liu_review_2011, sheridan_self-organizing_2011, hewitson_self-organizing_2002, gibson_use_2017}.
Regardless of the exact application, the trained SOM and its associated 2D node space must be interpretable by the user to be useful as this serves as the backbone for its use.
%
Past works have mostly considered Sammon Mapping or Multidimensional Scaling (MDS) to understand the SOM-defined 2D space; however, these methods do not provide user control over the mapping process.

% 
In \ClimateSOM, we extend SOM-based approaches by integrating interactive visual analytics techniques. In particular, we augment the SOM node space with interpretability tools including user annotations, interactive dimensionality reduction, and LLM-empowered bidirectional queries. This allows domain scientists to both guide and interrogate the structure learned by the SOM, moving beyond static representations and improving the transparency and utility of the latent 2D embedding.

%A recurring theme across these studies is use of SOMs as a clustering technique for organizing raw climate data.
%%
%% In fact, many works highlight its relationship to k-means clustering \cite{sheridan_self-organizing_2011, gibson_use_2017, hewitson_self-organizing_2002} and its capacity to extract realistic features from the input data.
%%
%Gibson et al. refers to this flavor of SOMs as \textit{SOM for clustering} \cite{gibson_use_2017}.
%%
%In these works, the relative importance is given to the expression in \textit{each} of the extracted SOM nodes since they can represent extreme events \cite{grotjahn_north_2016} or characteristic patterns
%of historical weather or climate variation within a region of interest \cite{loikith_characterizing_2017}.
%%
%This contrasts with \textit{SOM for projection} where the primary goal to uncover the underlying structure in the data.
%%
%Broadly, this application of SOMs can be considered as a form of dimensionality reduction.
%%
%% Related methods for dimensionality reduction in climate science including Empirical Orthogonal Functions (EOF) and Matrix Factorization (MF).
%%
%% However, these methods have its limitations including the interpretability of the extracted patterns due to their orthogonality constraints in EOFs \cite{liu_performance_2006} and the necessary restriction that they the archetypes in MF-powered methods lie on the extremes of the data \cite{harries_applications_2020}.
%%
%The classification of \textit{SOMs for projection} matches several works as well.
%%
%For instance, Radi\'{c} et al. used SOMs to analyze CMIP5 data over British Columbia \cite{radic_future_2015}, and more recently, Taylor et al. used SOMs to analyze CMIP6 data to identify recurring patterns in the CMIP6 ensemble \cite{taylor_cmip6_2023}.
%
% \ClimateSOM also follows this approach, using SOMs as a projection step to simplify the input ensemble data, reducing complexity and filtering out random variations.
%
% In \ClimateSOM, SOMs are adopted primarily as a processing step, emphasizing the characteristic of a set of nodes rather than the characteristics of individual nodes.
%%
%\ClimateSOM also aligns more with this use of SOMs, using SOM as a projection step for the input ensemble data, reducing the complexity and random variations the input ensemble data.
%%
%With \ClimateSOM, SOMs as adopted as processing step, emphasizing the \textit{distribution} of expressions in each of the input data, rather than the expressions of individual nodes. 

% \subsection{SOMs in Visual Analytics}

% SOMs have been used in a variety of visualization works in the past across different domains.

%Chen et al. used SOMs to visualize a collection of multi-objective pareto frontiers in project portfolio management (PPM) and Asset Management \cite{chen_self-organizing_2013}.
%%
%Hemmati et al. used SOMs to visualize and approximate the spectral energy distribution data for galaxies \cite{hemmati_bringing_2019}.
%%
%Sacha et al. used SOMs as a part of their VA system for time series data \cite{sacha_somflow_2018}.
%%
%Guo et al. used SOMs to visualize a time series of high-dimensional sales data \cite{guo_visualization_2006}, and Chusig-Muzo et al. used SOMs to cluster and visualize electronic health records into ``prototype patients'' \cite{chushig-muzo_data-driven_2020}.
%%
%More recently, Tu et al. used SOMs to layout extracted phrases from a text corpus into 2D for user exploration \cite{tu_phrasemap_2024} and Eirich et al. used SOMs to cluster and visualize data from electric engine acoustic signals \cite{eirich_irvine_2022}.
%
% Beyond these applications, SOMs have found substantial use in the analysis of cytometry data \cite{van_gassen_flowsom_2015, quintelier_analyzing_2021}.
%
% SOMs are used to visualize complex flow and mass cytometry data, and have been shown to be effective in identifying cell populations and their relationships.
%
% Van Gassen et al. introduced FlowSOM, a SOM-based visualization interface for cytometry data, where a minimum spanning tree was used to visualize the relationship between the extracted SOM nodes \cite{van_gassen_flowsom_2015}.


%\subsection{LLMs for Sensemaking in Visualizations}
%Past works identify that additive annotations can help with comprehrension and recall of visualizations, aiding in its sensemaking process~\cite{zhao_annotation_2017, rahman_qualitative_2025}.
%%
%Additionally, allowing users to express \textit{their} information needs by labeling specific visual elements can support knowledge generation~\cite{bernard_comparing_2018}.
%% Furthermore, providing users with the ability to express \textit{their} information needs via labeling specific visual elements can support the users' knowledge generation process \cite{bernard_comparing_2018}.
%% Supplying users with the ability to express \textit{their} information needs via labeling specific visual elements can support the users' knowledge generation process \cite{bernard_comparing_2018}.
%%Reasoning with complex data often requires its users to explore nuanced information spaces.
%%
%Meanwhile, recent studies highlight the potential of LLMs in facilitating sensemaking within visual analytics~\cite{hutchinson_llm_2024}.
%% Meanwhile, recent works demonstrate the potential for LLMs to aid in the sensemaking process of complex information for visual analytics \cite{hutchinson_llm_2024}.
%%
%For instance, LLMs have been used to assist in reasoning of spatiotemporal sensor traces~\cite{ouyang2024llmsense} and they have been also used to explore the latent space of text embedding models~\cite{raval_explain-and-test_2023}.
%%
%They have also been used to help users make sense of and explore the decomposition of information-seeking tasks~\cite{suh_sensecape_2023}.
%%
%They have also been shown to alleviate the ``cold-start problem'' in information-seeking, helping users establish context and find a starting point when navigating new information~\cite{liu_selenite_2024}.
%% whereby LLMs can provide a starting point and neccessary context for users as they explore new information \cite{liu_selenite_2024}.
%% Another recent work also demonstrates LLM's capability to help users to make sense of and explore complex information spaces like the decomposition of information-seeking tasks \cite{suh_sensecape_2023}.
%%
%% The premise of these works is that LLMs, given proper prompting and guardrails, can often prime users to new information spaces and explain patterns in data with minimal user intervention.
%The underlying premise of these works is that, with proper prompting and guardrails, LLMs can prime users to new information spaces and help explain patterns in data with minimal intervention.
%%
%With \ClimateSOM, we build on these insights by enabling users to provide \textit{annotations} within the SOM-defined 2D space, using LLMs in the process.

\subsubsection{LLMs for VIS}
Natural Language Interfaces (NLI)~\cite{shen_towards_2023} or LLMs have been increasingly applied to assist in visualization workflows~\cite{hutchinson_llm_2024, zhang_natural_2024, zhao_leva_2025}.
%
Such systems show promise in helping users by generating interpretable annotations in visualizations or by providing a conversational interface to the visualization system~\cite{wu_chartinsights_2024,shen_towards_2023} which can aid in their sensemaking process~\cite{bernard_comparing_2018,rahman_qualitative_2025}.
%
They have also been shown to alleviate the ``cold-start problem'' in information-seeking, helping users establish context and find a starting point when navigating new information~\cite{liu_selenite_2024}.
%
LLMs have also been successfully applied for Text-to-SQL tasks ~\cite{zhang_natural_2024, li_dawn_2024} and for summarizing insights from and improve usability in visual analytics interfaces~\cite{zhao_leva_2025}.
%
In fact, LLMs have been used in geospatial visual analytics systems as in \cite{li_save_2024} for gathering and summarizing insights in a geospatial context.
% or for generating SQL queries from natural language questions
%

In \ClimateSOM, we integrate LLMs more tightly into the system design by enabling bidirectional interaction between users and the SOM-defined latent space. This includes support for text-based queries into the embedding and automated summarization of regions in the latent space. This approach not only improves usability but also enhances the interpretive bandwidth of the SOM, particularly in settings where exploratory context is limited or evolving.



% With \ClimateSOM, we enable users to provide \textit{annotations} in the SOM-defined 2D space as a means to express their information needs, leveraging LLMs in the process.
%
% Meanwhile, the latent space of dimensionality reduction techniques have been freqeuntly explored in the past \cite{sacha_visual_2017}, many visual analytics works seeking to improve its interpretablity \cite{montambault_dimbridge_2025, eckelt_visual_2023}.
%
% Despite the documented issues like hallucinations that LLMs face, in light of these works, we propose the use of LLMs to help users navigate the latent space of the SOM in the \ClimateSOM workflow and reduce the "cold start" problem in exploring this latent space.
%Visualizing complex data often involves  
%dimensionality reduction to a lower-dimensional space, usually in 2D, to facilitate its understanding.
%%
%The premise is that patterns in the data are more easily discernible after processing; however, the problem of understanding the patterns still remains albeit in a lower-dimensional space.
%%
%Recent works have explored the use of LLMs to assist in sensemaking in these latent spaces \cite{raval_explain-and-test_2023}.


\subsection{Problem Characterization}\label{sec:problem}
We focus on ensemble datasets such as those produced by the Coupled Model Intercomparison Project (CMIP6)~\cite{eyring2016overview}, which contain climate projections generated under a variety of input assumptions. Each simulation produces a spatiotemporal time series, referred to as a model run. The following terminology is used throughout:
\begin{itemize}[noitemsep, topsep=0pt]
    \item \textbf{GCM} (Global Climate Model): Mathematical models that simulate the Earth's climate system, and are used to project future climate scenarios (e.g., MIROC6, ACCESS-CM2, etc.).
    \item \textbf{SSP} (Socioeconomic Pathways)~\cite{riahi2017shared}: Pathways that examine how global society, demographics, and economics might change over the next century. 
    In addition to the historical period, we consider three SSPs in this work: SSP245, SSP370, and SSP585, a low, medium high, and high greenhouse gas emissions future respectively. 
    \item \textbf{Model run}: A unique climate projection and spatiotemporal time series, representing a member of the climate ensemble generated by a specific combination (GCM, SSP). The full ensemble consists of multiple such runs.
\end{itemize}

\subsubsection{Design Goals}\label{ssec:design_requirements}
Throughout the project, in consultation with a climate scientist (paper co-author) specializing in regional climate variability and changes, we identified three key questions that drive exploratory analysis of ensemble climate datasets. These questions inform the system's design goals.
% Given a climate ensemble dataset, we list the motivating questions and the design requirements for \ClimateSOM.
\begin{tightItemize}
    \item[\req{R1}] {\normalfont \sffamily``Q: What is the general behavior of the set A of model runs?''}  \\ A climate scientist wishes to understand the general behavior of a single model run or a set of model runs and its general tendencies as it relates to the climate projection it generates. For example, do the model run(s) consistently produce projection of a certain quality or does it vary over time? \\ $\rightarrow$ \textbf{Requirement: }\textit{Provide a visual language and description to explore the behavior of a single set of model runs over the entire temporal domain.}
    \item[\req{R2}] {\normalfont \sffamily``Q: How does set A of model runs compare to set B of model runs?''} \\ A climate scientist wishes to compare the behavior of two sets of model runs A and B to understand the differences in their projections. For example, does set A of model runs produce similar projections compared to set B of model runs? What overlapping qualities do they share and what climate projection qualities are unique to run A or run B?\\ $\rightarrow$ 
    \textbf{Requirement: }\textit{Provide a visual language description and support to compare the behavior of two sets of model runs over the entire temporal domain.}
    \item[\req{R3}]{\normalfont \sffamily``Q: Which model runs or GCMs exhibit \textit{similar} behavior?''} \\ A climate scientist wishes to understand whether there are meaningful clusterings of ensemble runs and clustering of GCMs based on the SSP forcings. Do all ensemble runs produce similar results? Is the effect of SSP forcing similar on all GCMs?\\ $\rightarrow$ 
    \textbf{Requirement: }\textit{Support clustering of ensemble model runs based and clustering of GCM based on their SSP forcings over the entire temporal domain.} 
\end{tightItemize}
% The goal of \ClimateSOM is to provide answers to the above question while ensuring that they are readily decipherable.
% The goal of \ClimateSOM is to address the above questions with an interpretable visualization tool.
% \note{How to properly incorporate interpretablitity as a design requirement?}
% As outlined in~\cref{sec:intro} the main task

\subsubsection{Abstracting the Model Runs Into a Distribution}\label{ssec:abstraction}
Central to \ClimateSOM is a self-organizing map (SOM) trained to capture the spatial variability of a climate ensemble dataset, while retaining its interpretability.
%
An SOM is an unsupervised neural network that projects high-dimensional data into a 2D grid while preserving the topological structure, making it a useful tool dimensionality reduction (DR) of its input.
%
Training a SOM over a climate ensemble dataset yields a 2D SOM grid, each SOM node representing a spatial pattern over the studied region.
%
\ClimateSOM leverages this capability to represent each model run within the ensemble as a distribution.
%
Specifically, for each time step we first locate its best matching unit (BMU) in the SOM nodes, and allow the 2D position of that SOM node to represent that time step.
%
Aggregating these positions across all time steps produces a 2D distribution defined by the set of BMU's node locations.
%
\cref{fig:SOM} presents the construction of the distributional representation.

Noteworthy in this construction, however, is the projection step of the initial 2D SOM grid. 
%
While the SOM nodes maintain local smoothness, they can have distortions globally~\cite{sheridan_self-organizing_2011}; thus, defining the SOM node locations solely by their 2D grid location is not desirable.
%
To address this, we add a projection step of the SOM nodes to better preserve their relative structure and enhance its interpretability.
% 
%We thus aim to correct for these distortions by projecting the SOM nodes into a more interpretable orientation to better preserve the relative structure of the SOM nodes.
%
% Finally, we use the projected SOM node locations to construct the distributional representation as shown in~\cref{fig:SOM}.
\begin{figure}[H]
   \includegraphics[width=\linewidth]{SOM-10.pdf}
    \caption{The distributional abstraction for model runs in \ClimateSOM. Each model runs are abstracted to a distribution defined by a projected location of SOM Nodes, the \textit{SOM Node Space}.}
    \label{fig:SOM}
\end{figure}
\noindent
Formally, for an ensemble run with $n$ time steps, each a climate projection over some spatial domain, $R=[r_1, r_2, \ldots, r_n]$, we abstract this as a set of points in 2D space:
\begin{align*}
S=\{p_1,p_2,\ldots,p_n\}, \,\,\, p_i=\text{proj}(BMU(r_i))
\end{align*}
where $BMU(r_i)$ is the best matching SOM node for $r_i$, and proj$(\centerdot)$ maps the node to a 2D layout.
%
This distributional representation of the model runs serves as the backbone of \ClimateSOM.
%
% Therefore, enabling the interpretation of the projected space of the SOM nodes is principal to our workflow which we


% as a set of points $S=\{p_1, p_2, \ldots, p_n\}$ in an interpretable 2D space, where each $p_i$ is the dimensionality-reduced position of best matching unit (BMU) of $r_i$ within the trained SOM.
%
%Designed to achieve these qualities, the workflow of \ClimateSOM is shown in~\cref{fig:workflow} and is composed of the following steps:
\subsubsection{Workflow Overview}\label{ssec:workflow_overview}
The workflow contains one pre-step and three main steps (\cref{fig:workflow}).
\begin{itemize}[left=0pt, itemsep=0pt, labelsep=1em, align=left]
    \item[\step{S0} \textbf{Training the SOM}]
        The workflow begins with training the SOM on the dataset to ensure it has sufficient explanatory power---specifically, that the approximation via its BMU is adequate---and smoothness across the map.
        % A crucial assumption of \ClimateSOM is that SOM must capture \textit{sufficient} variance in the ensemble data, so that $r_i \approx BMU(r_i)$ is a reasonable approximation.
        % Careful consideration is required for the parameters for the SOM training and the metrics with which to evaluate and these are described in~\cref{ssec:som_training_parameters,ssec:som_training_metrics}.
        % The SOM is trained to capture the spatial variability of the input data set. The SOM is trained using the \textit{Kohonen} algorithm~\cite{kohonen1982self}, which is a type of unsupervised learning algorithm that is used to produce a low-dimensional representation of the input data set.
    \item[\step{S1} \textbf{Anchor}]
        Using the pseudo-2D smooth SOM nodes from Step 0, this step enables users to create the Adjusted SOM Node Space by projecting the SOM nodes into a 2D layout. To support interpretation and navigation, users can anchor specific nodes---fixing their positions while allowing others to adjust---resulting in a more interpretable spatial arrangement than the original SOM grid. 
        We refer to this refined space as the \textit{Adjusted SOM Node Space}.
        %
        % Here, instead of using common techniques like t-SNE\cite{van2008visualizing} or UMAP\cite{mcinnes2018umap}, we opt for Minimum-Distortion Embedding (MDE) \cite{agrawal2021minimum}, which generalizes t-SNE and UMAP and crucially allows the anchoring of points in the output 2D space.
        % %
        % When supplied with a anchoring constraint, MDE ensures that these points are \textit{fixed} in a given location in the output space, while reducing the distortion of the remaining points.
        % %
        % We prefer MDE over t-SNE and UMAP since often fail to preserve important characteristics of the embedded space like interpretable axes.
        %
        % We use MDE to allows users to iteratively construct the \textit{Adjusted SOM Node Space}, anchoring points within the embedded 2D space to ensure a more meaningful layout of the SOM nodes.
    \item[\step{S2} \textbf{Annotate}]
        Once the \textit{Adjusted SOM Node} Space is defined, users annotate it to produce the \textit{Annotated SOM Node Space}. An annotation consists of (1) a bounded region within the adjusted space and (2) a text label describing that region. These annotations serve as an abstraction layer to facilitate interpretation in the next step. \ClimateSOM incorporates LLMs to assist with this process.
    \item[\step{S3} \textbf{Analyze}]
        A given model run can now be described as a distribution within the \textit{Annotated SOM Node Space}, as well as percentages broken down by user-defined \textit{annotations} \req{R1}. Model runs can be compared either by directly comparing their distributions or by interpreting them as vector fields over the annotated space \req{R2}. Clustering can be applied to model runs based on these distributions, and to SSP forcings over each GCM by clustering changes in their respective distributions \req{R3}. Please see the details in~\cref{sec:workflow}.
\end{itemize}

\begin{figure*}[h]
    \centering
    \includegraphics[width=\textwidth]{workflow.pdf}
    \caption{The \ClimateSOM workflow. The workflow is composed of the SOM Training step and three main steps: Anchor, Annotate, and Analyze. In \step{S1}: Anchor, the user adjusts the SOM Node Space using iterative anchoring with Minimum-Distortion Embedding to produce the \textit{Adjusted SOM Node Space}(\subimage{2}). Dragged nodes are called \textit{anchors} and are fixed in place, while the rest of the nodes are adjusted. In \step{S2}: Annotate, the user annotates the SOM Node Space using manual inspection of the SOM Nodes or with the LLM Integration. Once annotated, in \step{S3} the user can
        use the \textit{Annotated SOM Node Space} (\subimage{3})
    as the basis to analyze the ensemble dataset.}
    \label{fig:workflow}
\end{figure*}



\noindent
Having outlined the workflow, we now address two important design choices: the use of SOMs and the intentional abstraction away from temporal ordering.

\vspace{3pt}
\noindent\textbf{SOM as the central tool.} 
We choose SOMs over other DR or clustering methods for the following two reasons:
\begin{enumerate}[noitemsep, topsep=0pt]
\item \textbf{Preservation of local topology}: SOMs preserve neighborhood relationships in lower-dimensions which enables meaningful spatial interpretations. Our distributional abstraction of model runs requires that the learned grid be continuous and neighboring nodes express similar spatial patterns. Other DR or clustering methods do not produce topologically smooth results in lower-dimensions and thus do not satisfy the conditions required in \ClimateSOM.
\item \textbf{Established use in climate science}: SOMs are widely used for analysis of climate data and are familiar to domain experts, aiding in its interpretation.
% \item \textbf{Visual continuity and smoothness}: Our data distributional abstraction of model runs requires that the learned grid must be continuous and neigh
\end{enumerate}

% Among the various possible methods to initially reduce the dimensionality of the input ensemble, \ClimateSOM opts for SOM for two main reasons: (1) SOMs are an established tool in climate science and are known to be successful in extracting salient features from spatiotemporal climate data and (2) \ClimateSOM relies on preserving the variability of the input ensemble via a pseudo-smooth 2D representation.  
%
% Maintaining smoothness in this representation is crucial, as discontinuities or irregularities can compromise the interpretability of the abstraction.
%
% While smoothness can be defined in various ways, we require that, locally at each node, the neighborhood maintains a consistent level of smoothness across the grid. 
%
% This ensures that the MDE embedding, powered by local neighborhood graphs, produces interpretable axes in 2D.
% \note{How do I properly define meaningful after MDE?}
%
% SOMs with the tuned parameters can guarantee this; however, other methods like hierarchical or \textit{k}-means clustering do not.

\vspace{3pt}
\noindent\textbf{Ignoring the temporal ordering.}
\ClimateSOM represents a model run $R$ as a set of points $S$, thus discarding temporal order (scrambling R yields the same S). This is an intentional trade-off for two reasons:
\begin{enumerate}[noitemsep, topsep=0pt]
    \item \textbf{Emphasis of broad patterns}: The goal is to analyze the general behavior of model runs over time rather than their specific sequence. This align with how GCM outputs are studied for broad differences in climate science.
    \item \textbf{Reducing complexity}: While preserving temporal order enables the use of time series analysis techniques, such as dynamic time warping, it also introduces complexity and amplifies the random variations inherent in climate ensemble data.
\end{enumerate}
This trade-off enables simpler interpretation while retaining the ability to extract broad patterns from model outputs.
% %
% (1) \
% %
% (2) while accounting for the temporal order allows for the application of time series  analysis and techniques like dynamic time warping, it is certain that it also introduces complexity and accentuates the random variations that climate ensemble data exhibits.
%
% Despite the loss of the temporal ordering, we argue that \ClimateSOM can still produce meaningful insights into the behavior of climate ensemble datasets.

\subsection{ClimateSOM}\label{sec:workflow}
% We describe the steps of the \ClimateSOM workflow in detail in this section.
\subsubsection{Training the SOM (\step{S0})}
% \subsection{Self-Organizing Maps (SOM)}
% Self-organizing maps are a type of neural network designed to learn the topology of the input data by mapping it to a lower-dimensional grid, usually in a rectangular grid in 2D of size $dim\times dim$.
%
% We refer to \cite{kohonen_self-organizing_1990} for its formal definition.
%
% After a random initialization, at each iteration, the SOM algorithm selects a random input field and finds the best matching node (BMU) in the grid. 
%
% While varying definitions of \textit{best matching} exist, we use the Euclidean distance between the input field and the node's field as the distance metric.
%
% After identifying the BMU at each training step, the algorithm updates the SOM weights based on a learning rate and a neighborhood function.
%
% Formally, for a data set $R=\{r_1, r_2, \dots r_n\}$ of a collection of spatiotemporal fields (in the case studies, the fields are a series of precipitation anomalies over chosen domains in the Western US), at every iteration $t$, the algorithm selects a random sample $x(t)$ from $R$ finds its BMU, node $b$, and the weight vector $w_i$ of the node $i$ is updated as follows:
% \begin{equation}\label{eq:som_update_rule}
%     w_{i}(t+1) = w_{i}(t) + \alpha(t)h_{i,b}(t)[x(t) - w_{i}(t)]
% \end{equation}
% where $w_{i}$ is the weight vector of the node $i$, $\alpha(t)$ is the learning rate at time $t$, $h_{i,b}(t)$ is the neighborhood function that computes the influence of node $i$'s update, defined as a function of the distance between nodes $i$ and $b$.
%
For our spatiotemporal ensemble dataset, we train the SOM on the flattened ensemble dataset (i.e. across all spatial realizations that occur in the ensemble).
%
As common choice, we opt for the neighborhood function $h(t)$ to be the Gaussian function and linearly decrease its $\sigma$ as is standard practice in climate applications~\cite{gibson_use_2017}.
%
After training, the grid of SOM nodes is expected to approximate and learn the distribution input ensemble dataset under the constraint that the 2D grid of nodes are pseudo-smooth.

The SOM must meet two competing criteria: (1) it should retain sufficient information from the input data, and (2) the 2D grid of nodes should be smooth enough to avoid disconnections after projection, which would hinder interpretability.
%
However, these two criteria are often in conflict. 
%
% Smaller neighborhood update size often results in higher explanatory power allowing the SOM to extract finer details of the input data's topology; however, this is at the cost of non-smooth SOM grid since each update influences a smaller region of the grid, and vice versa.
%
% A smaller neighborhood update size enhances the SOM’s ability to capture finer details of the input data’s topology, increasing its explanatory power. 
%
% However, this comes at the cost of a less smooth SOM grid, as each update affects a smaller region. 
%
% Conversely, a larger neighborhood size promotes smoother transitions but may obscure finer details.
%
A smaller neighborhood update size captures finer details of the data’s topology but results in a less smooth SOM grid, as each update affects a smaller region. 
%
Conversely, a larger update size smooths the grid but may obscure finer details.
%
We found that the choice of the initial neighborhood update size ratio expressed as a ratio between the initial update size and the dimension of the SOM grid, ($kR=(\sigma_{initial})^2/ {dim}^2$), and the ratio of the final neighborhood update size, $kS=\sigma_{final}/ \sigma_{initial}$, to the initial size ($kS$) are critical to balancing these two criteria.
%
For our case studies, we found that $kR=0.03$ and $kS=0.2$ provide balanced results.
%
Further details are in the Appendix.



%\subsubsection{On the metrics for SOM training}\label{ssec:som_training_metrics}
%Past research have employed various metrics to evaluate the quality of a SOM \cite{forest2020survey}.
%%
%Generally, metrics for SOM are divided into those that evaluate error that the SOM approximation introduces and those that evaluate the topological properties of the SOM grid.
%%
%\ClimateSOM requires both classes of metrics to measure that the trained SOM abstraction retains \textit{sufficient} information whilst also ensuring that the MDE dimensionality-reduced space is connected and interpretable.
%%
%The most common metrics for the respective categories are quantization error ($QE$) and topographic error ($TE$).
%%
%% For a given dataset $D=\{d_1, d_2, \ldots, d_n\}$, the quantization error is defined as the average distance between each input field $d_i$ and its BMU, $BMU(d_i)$.
%%
%% Topographic error is defined as the proportion of all data for which the first and second BMUs are \textit{not} adjacent in the SOM grid.
%%
%% While these are useful metrics, we introduce two more metrics that we have found useful for \ClimateSOM: (1) Explained Variance ($EV$) and (2) Variation of Smoothness ($VS$).
%While these are useful metrics, we present two more metrics that we have found useful for \ClimateSOM: (1) Explained Variance ($EV$) and (2) Mean Smoothness ($\mu_S$).

%\noindent
%\textbf{Explained Variance}, akin to its definition in principal components analysis (PCA), measures the proportion of variance in the input data set that is explained by the SOM grid, where 0 indicates no explanatory power and 1 indicates full explanatory power.
%%%
%%Formally, for a data set $D$ the explained variance is defined for a trained SOM is defined:
%%\begin{equation}
%%    EV = 1 - \frac{\sum_{i=1}^{n}\sum_{j=1}^{m}(d_{i_j} - (BMU(d_i))_j)^{2}}{\sum_{i=1}^{n}\sum_{j=1}^{m}(d_{i_j} - (\bar{d})_j)^{2}}
%%\end{equation}
%%%
%%where $n$ is the number of members in $D$, $m$ is the number of spatial locations, $d_{i_j}$ is the $j^{th}$ element of the $i^{th}$ input field, $BMU(d_i)_j$ is the $j^{th}$ element of the BMU of the $i^{th}$ input field, and $\bar{d}_j$ is the $j^{th}$ element of the mean of the input data set.
%%% 
%We replace $QE$ with Explained Variance to provide a unitless measure of the explanatory power of the SOM grid.

%%\textbf{Variation of Smoothness} is another metric that we introduce to more carefully evaluate the smoothness of the SOM grid.
%%%
%%Given a SOM grid $G$, we define local smoothness, $LS$, at any grid node $g_{i,j}$ as the average Euclidean distance between $g_{i,j}$ and its four immediate neighbors.
%%%
%%Finally, considering the standard deviation of all $g_{i,j}$, we define the Variation of Smoothness as $\sigma_{LS}/ \mu_{LS}$.
%%%
%%Lower values of Variation of Smoothness indicate a more consistent level of smoothness across the grid and vice versa.
%%%
%%We found that topographic error was not sufficient to capture the required smoothness of the SOM grid to ensure the interpretablitity of the MDE result and thus replace the topology-based error metric with the Variation of Smoothness.

%%In all, SOM training for \ClimateSOM is the balancing of its explanatory power ($QE,EV$) and smoothness ($VS$).
%%%
%%The following, \cref{ssec:som_training_parameters}, discusses the choice of parameters in the context of these metrics and the implications.
%\noindent
%\textbf{Mean Smoothness} is another metric that we introduce to more carefully evaluate the smoothness of the SOM grid as non smoothness can produce disconnected MDE results.
%%
%Given a SOM grid $G$, we define local smoothness, $LS$, at any grid node $g_{i,j}$ as the average Euclidean distance between $g_{i,j}$ and its four immediate neighbors and define the Mean Smoothness as $\mu_{g_{i,j}}$.
%%
%We found that topographic error was not sufficient to capture the required smoothness of the SOM grid to ensure the interpretablitity of the MDE result and thus replace the topology-based error metric with Mean Smoothness.

%In all, SOM training for \ClimateSOM is the balancing of its explanatory power ($QE,EV$) and smoothness ($\mu_S$).
%%
%The following, \cref{ssec:som_training_parameters}, discusses the choice of parameters in the context of these metrics and the implications.


%\subsubsection{On the parameter choices for SOM training}\label{ssec:som_training_parameters}
%Among the parameters required for SOM training the initial neighborhood size $(kR)$ and the neighborhood size decay $(kS)$ are most important and we thus list their influences here.
%\begin{itemize}[left=5pt, itemsep=0pt, labelsep=1em, align=left]
%    % \item[\textbf{dim}] The dimension of the SOM grid defines the number of nodes \ClimateSOM uses to approximate the variations in the input ensemble. Larger values tend to result in better explanatory power, but adds complexity to the MDE. While non-square grids can be used, we opt for square grids for simplicity.
%    \item[\textbf{kR}] The Gaussian neighborhood function depends on the parameter $\sigma$. The term $kR$ defines its initial value as the ratio between the SOM grid area and the neighborhood update area, given by $kR=(\sigma_{initial})^2/ dim^2$. 
%        %
%        % Generally, a larger $kR$ results in a smoother SOM grid with less explanatory power (i.e. low $VS$, low $EV$ and high $QE$) and a smaller $kR$ results in the opposite (i.e. high $VS$, high $EV$ and low $QE$).
%        Generally, a larger $kR$ results in a smoother SOM grid with less explanatory power and a smaller $kR$ results in the opposite.
%    \item[\textbf{kS}] In order for the SOM to \textit{fit} the data, it is common to linearly decay the neighborhood size over time \cite{gibson_use_2017}. %
%        $kS$ defines the final value of the neighborhood size as a ratio of the initial value, $kS=\sigma_{final}/ \sigma_{initial}$.
%        %
%        Similar to $kR$, $kS$ also expresses a trade-off between smoothness and explanatory power.
%        %
%        A larger $kS$ results in a smoother SOM grid with less explanatory power and a smaller $kS$ results in the opposite.
%\end{itemize}
%Other parameters such as the dimensions of the SOM grid, as long as they were sufficiently large to enable sufficient explained variance, learning rate decay, the number of iterations, and the initialization method are part of the training process but were found less critical to the performance, and hence we omit them from discussion.

% \begin{figure*}[h]
%     \centering
%     \includegraphics[width=0.9\textwidth]{figures/TaskDescriptions.pdf}
%     \caption{The tasks that \ClimateSOM accomplishes and its overview}
%     \label{fig:tasks}
% \end{figure*}
% \begin{figure}[t]
%     \centering
%     \includegraphics[width=0.9\linewidth]{figures/LLM.pdf}
%     \caption{Incorporatings LLMs into the \ClimateSOM workflow}
%     \label{fig:llm}
% \end{figure}
% \begin{figure}[t]
%     \centering
%     \includegraphics[width=0.9\linewidth]{figures/BootstrapComparison.pdf}
%     \caption{The bootstrap-based construction of comparison within the \ClimateSOM workflow}
%     \label{fig:bootstrap_comparison}
% \end{figure}


\subsubsection{Anchor (\step{S1})}
% \subsubsection{Dimensionality Reduction (DR) on the SOM Nodes}
As seen in past works in climate science employing SOMs \cite{sheridan_self-organizing_2011}, the SOM grid, while constructed to be locally smooth due to the training process, can have distortions globally. 
%
Therefore, SOMs are often presented alongside Sammon maps or Multidimensional Scaling (MDS) to provide a more faithful view of the SOM grid in terms of its relative distances between nodes.
%
However, these methods give no control to users to control the mapping process and gives priority to global structure in place of local structure.
%
To provide users with the iterative ability to shape and project the SOM nodes, we allow user to \textit{anchors} particular SOM node in the projection process.
% 
% In line with these approaches, we propose to apply a dimensionality reduction (DR) technique on the SOM grid nodes to yield a more interpretable 2D representation of the spatiotemporal fields in the input data set.
%
% However unlike similar approaches that rely of global DR methods like Sammon mapping or MDS, we propose a local DR method for two reasons. 

% First, the most interesting comparison arises between ensemble members whose variations are often local in nature. 
%
% First, the differences between ensemble members can be subtle since they all model similar phenomena and subtle differences are better captured by local DR methods.
%
% In downstream tasks like clustering or the characterization of forcing, we rely on collection of local differences to encapsulate overall differences, and thus preserving local structures is preferred over its counterpart.
%
% Second, Euclidean distance, our distance metric to define pairwise difference between SOM nodes, is a sufficient local metric but a poor global metric, and thus optimizing for the global structure may introduce distortions that are less desirable.
%
% An important component of this step is the anchoring of points in the output 2D space as a constraint to the DR process.
%
% This allows for the user to configure the space in which the ensemble data will be visualized by explicitly soliciting the user to reorganize the 2D space to better suit their needs.
%
The utility of placing constraints on DR processes has been explored in previous works \cite{vu_integrating_2022}, and we find the iterative nature of anchoring via to be most suitable as it requires minimal intervention from the user.
%
% We claim that this is preferable over methods like UMAP as these methods can fail to preseve important characteristics of the embedded low-dimensional space, such as the smooth contours of the SOM grid.
%
In \ClimateSOM, we employ Minimum-Distortion Embedding (MDE)\cite{agrawal2021minimum} to support anchoring.
%
Similar to UMAP\cite{mcinnes2018umap}, MDE can be presented as an optimization problem over local neighborhood graphs.
% , where the distortion between any two input data $i$ and $j$ arising from the dimensionality reduction is defined as a function of $d_{i,j}$, the Euclidean distance between the $x_i$ and $x_j$ -- embeddings of $i$ and $j$. 
%
%This function is formalized by Agrawal et al. as $f_{(i,j)}(d_{i,j})$ and referred to as the \textit{distortion function}.
%%
% Letting $\varepsilon$ to be the set of all pairs of $(i,j)$ and defining the \textit{average distortion} as $E(X)=(\Sigma_{(i,j)\in\varepsilon}f(i,j)) / |\varepsilon|$, MDE is formulated as the optimization problem to minimize $E(X)$ over all sets of feasible embeddings of the input $X$.
%
% Depending on MDE's configuration, this generalization can represent various DR techniques from PCA, MDS to UMAP.
%
Using MDE, we optimize for the distortion of the SOM grid nodes in the 2D space while anchoring points in the output space to ensure a more meaningful layout of the SOM nodes.
%
Since a 2D SOM roughly extracts two main axes of variation, this step allows user to physically map these axes to the two axes of the visual interface itself.
%
We refer to the output of this step as the \textit{Adjusted SOM Node Space}.

\subsubsection{Annotate (\step{S2})}\label{ssec:annotate}
This step allows users to provide \textit{annotations} to the \textit{Adjusted SOM Node Space} to add user-defined context to the SOM nodes while preparing for subsequent analysis in \step{S3}.
%
We define an \textit{annotation} as a pair of the following: (1) a bounded region within the \textit{Adjusted SOM Node Space} and (2) a text label that describes the region.
%
Two challenges arise: identifying relevant regions and defining meaningful labels.
%
Given the documented success of Text-to-SQL via LLMs \cite{zhang_natural_2024} and text summarization \cite{pu_summarization_2023}, our proposal for \ClimateSOM is leveraging LLMs in this process, alongside the user's visual inspection of the individual SOM nodes.
%
These correspond to the forward and backward directions of LLM assistance as shown in \cref{fig:workflow}.


\paragraph{Forward Direction}
In the forward direction, we assume the user has a characteristic of interest in mind for SOM nodes (e.g., ``\textit{Show me SOM nodes with average precipitation over Southern California above 0}''). The goal is to identify the corresponding region in the \textit{Adjusted SOM Node Space} by filtering nodes that match the specified condition. This is done by generating a PostGIS query based on the user’s geospatial query.
%
However, we found that naively prompting this task to LLMs does not yield consistent results, primarily due to LLM's
% An identified challenge with tasking LLMs to generate spatial queries from natural language, however, is the 
lack of accuracy in generating appropriate boundaries for given geospatial regions as found in \cite{liu_nalspatial_2025}.
%
To address this, we delegate the LLM not to generate spatial boundaries, but instead to generate a list of counties that correspond to the named region. These counties, which have well-defined boundaries, are then used to construct the spatial filter, reducing errors introduced by the LLM.
%
% In addition, we also adopt self-correction for the SQL query generation as seen in \cite{pourreza_din-sql_2023}.
%
In all, the forward direction proceeds as follows.
\begin{enumerate}[nolistsep]
    \item Parse the user query to detect any region names or locations. If none are detected, skip to step 3.
    \item Resolve the identified regions to a list of counties with known boundary definitions. Combine these county boundaries to form the target region.
    \item Use the resolved region and the numerical conditions extracted from the query to construct and execute a PostGIS query over the SOM nodes.
\end{enumerate}
\paragraph{Supported Question Types}
Based on discussions with the domain experts, we identify three categories of user queries that the LLM should be able to support by generating appropriate PostGIS queries:
%
\begin{enumerate}[nolistsep]
    \item SOM nodes such that the average value over \textit{some geospatial region} is above / below \textit{X}
    \item SOM node such that the average value over \textit{some geospatial region} is between \textit{X} and \textit{Y}
    \item SOM node such that the average value over \textit{some geospatial region A} is higher or lower than the average value over \textit{some geospatial region B}
\end{enumerate}


\begin{figure*}[!t]
    \centering
    \includegraphics[width=\textwidth]{temp.pdf}
    \caption{
        The \ClimateSOM interface at the end of \step{S1} and \step{S2} displaying the \textit{Annotated SOM Node Space}.
        %
%        The visual interface to support the first two steps --- \step{S1}: Anchor and \step{S2}: Annotate. 
%        %
%        The assumed user is first presented with \highlight{1}. 
%        %
        The user can inspect each SOM node via clicking on the node \subimage{1}, and the Node Viewer on the left displays the selected map of the SOM Node.
        %
        % After some initial visual inspection, they realize that the initial embedding of the SOM Nodes are not optimal since they cannot well interpret the dimensionality reduced axes. 
        %
        % The user then drags SOM Nodes around \subimage{1b} and clicks \textit{Reflect Anchors} and iteratively arrives to \highlight{2}. 
        %
        To support \step{S2}: Annotate, users can leverage the two directions of LLM-powered query system: Forward \subimage{2} and Backward \subimage{3}. 
        %
        For the forward direction \subimage{2}, the user can request a particular quality of interest, (e.g., \textit{Show me nodes that have average above 0 over Southern California}) and be presented with the corresponding SOM Node region highlighted in yellow. 
        %
        For the backward direction \subimage{3}, the user can select a SOM Node region and a LLM will generate a \textasciitilde
 50-word summary of SOM Node realizations in that region.
        %
        % The user's interaction with the LLM are tracked in \subimage{2d}.
        %
        Through visual inspection and the LLM assistance, the user compiles the set of \textit{annotations} as regions in the SOM Node Space and a label corresponding to the region.}
        \vspace{-0.2in}
    \label{fig:Step1and2}
\end{figure*}

\paragraph{Backward Direction}
In the backward direction, we presume the user has identified a region within the \textit{Adjusted SOM Node Space} and wishes get a high-level summary of the region. 
%
In \ClimateSOM we accomplish this as follows.
\begin{enumerate}[nolistsep]
    \item Given a region in the \textit{Adjusted SOM Node Space}, we extract samples evenly within the SOM Node Space region. 
        % For our case studies, we found nine samples to yield sufficient results.
    \item Construct a per-county level summary statistic of each node from step 1. For a given sampled SOM node, we textually represent each node as: \begin{verbatim} 
{ high_value:[counties_list], 
    moderately_high_value:[counties_list],
    neutral_value:[counties_list],
    moderately_low_value:[counties_list],
    low_value:[counties_list] }
\end{verbatim} 
    That is, we enumerate all counties with average values in each of the five categories.
    %
    Note that, we allow the the cutoff values that define \textit{high}, \textit{moderately high}, etc. to be modified by the user.
    \item Use the text-based summary statistics for each sampled SOM node with a LLM to generate a summary description of the region.
\end{enumerate}

\vspace{3pt}
\noindent
In summary, the LLM is used in the forward direction for two tasks: (1) mapping region names to a list of counties, and (2) generating PostGIS queries to retrieve relevant SOM nodes. In the backward direction, the LLM serves as a summarizer of per-county statistics derived from user-selected regions.
%
We implement both directions using LangChain \cite{langchainLangChain} with OpenAI's \textit{gpt-4o-mini}\cite{achiam2023gpt} and use LangChain Tools for the forward direction to detect if the user input includes any geospatial regions that need to be resolved.
%
Details of the prompts used are available in the Appendix.
%
Either via manual inspection of the \textit{Adjusted SOM Node Space} or via the LLM integration, the user can now define a set of \textit{annotations} for the \textit{Adjusted SOM Node Space} to produce the \textit{Annotated SOM Node Space}.

\vspace{-3pt}
\subsubsection{Analyze (\step{S3})}
Finally in this step, we describe how \ClimateSOM performs each of the tasks and meets the design requirements outlined in~\cref{ssec:design_requirements}.
% \vspace{-3pt}
\paragraph{Exploring a single set of model runs \req{R1}}\label{ssec:explore}
% \begin{figure}[htbp]
%     \centering
%      \includegraphics[width=\linewidth]{figures/ExploreOne.pdf}
%      \caption{}
% \end{figure}
Given a single ensemble model run with $n$ time steps, $R=[r_1, r_2,\ldots, r_n]$, we can describe it as a distribution $S=\{p_1,p_2,\ldots, p_n\}$ within the \textit{Annotated SOM Node Space} by mapping each $r_i$ to its best BMU and applying the MDE projection: $p_i = MDE(BMU(r_i))$.
We then compute a kernel density estimate (KDE) over $S$, which serves as a “fingerprint” of the model run. Regions of high density indicate spatial patterns that occur frequently in the run. This distribution can also be expressed as percentages over the annotated regions defined in \step{S2}, enabling semantic interpretation beyond the KDE plot.
%
In the case of a set of multiple model runs, we apply the same process over the entire set of consider the aggregated set of points as the representative distribution.
\paragraph{Comparing the behavior of two sets of model runs \req{R2}}
% \begin{figure}[htbp]
%     \centering
%      \includegraphics[width=\linewidth]{figures/Compare.pdf}
%      \caption{}
% \end{figure}
To facilitate comparison of two model runs $R_1$ and $R_2$, \ClimateSOM provides two methods: (1) side-by-side comparison and (2) a vector field formulation.
%
\paragraph{Comparison via side-by-side}
The side-by-side comparison is a direct comparison of the distributions of $R_1$ and $R_2$ in the \textit{Annotated SOM Node Space} as defined in \cref{ssec:explore}.
%
By presenting these two views side-by-side, we allow users to visually inspect if there are regions within the SOM Node Space that are occupied in similar densities or different densities by $R_1$ and $R_2$.
%
Similar to the exploration of a single model run, the respective distributions can be compared via the breakdown based on the annotations defined in \step{S2} Annotate. 

\paragraph{Comparison via a vector field formulation}
We also support a formal comparison using vector fields. This is based on a bootstrap approach \cite{efron1992bootstrap} (implementation details in the Appendix).
%
Given two model runs $R_1$ and $R_2$, we generate $k$ number of bootstrap samples for each model run.
%
At each iteration, we first calculate the 2D optimal transport\cite{flamary2021pot} between the two bootstrap samples.
%
Given the optimal transport plan, we consider each pair in the plan and construct a set of vectors that points from each source to its target.
%
Interpolating the results for all sets of vectors onto a regular grid to generates a vector field.
%%
%Letting $R_1',R_2'$ be a bootstrap samples, for each bootstrap iteration we first compute the 2D optimal transport\cite{flamary2021pot} between $R_1'$ and $R_2'$.
%%
%This results in a list of pairs of elements $(r_1',r_2')$ where $r_1'\in R_1'$ and $r_2'\in R_2'$ such that the total cost of transport or matching is minimized.
%%
%Using this 2D optimal transport result, we aggregate and linearly interpolate this result onto a regular grid to generate a vector field.
%%
%Finally, averaging this result from all bootstrap iterations, we can generate a vector field that represents the difference between $R_1$ and $R_2$ in the \textit{Annotated SOM Node Space} as a transition from $R_1$ and $R_2$.
%
% Thus, the critical assumption of this approach is that is that the \textit{change} from $R_1$ to $R_2$ can be estimated by the 2D optimal transport between the two distributions.
%
% For our case studies, we found $k=50$ to produce stable results.
%
The resulting field represents the spatial transition from $R_1$ to $R_2$. 
%
This further enables analysis at the level of annotations---e.g., determining where points from annotation A in $R_1$ are mapped in $R_2$.

This approach is particularly useful when comparing temporally ordered scenarios (e.g. comparing a historical run to a future run or comparing a low to a high emission scenario). The vector field comparison may provide an added layer of context since the comparison is treated as a \textit{transition} from one to the other.
%
If comparing between two sets of model runs, we apply the same process to the aggregated set of points in each set as was in \req{R1}.



% \begin{figure}[h]
%     \centering
%     \includegraphics[width=\linewidth]{figures/BootstrapComparison.pdf}
%     \caption{The bootstrap-based construction of comparison for \ClimateSOM as a vector field}
%     \label{fig:bootstrap_comparison}
% \end{figure}

\vspace{-3pt}
\paragraph{Clustering the behavior of many model runs \req{R3} } \label{ssec:clustering}
% \setlength{\intextsep}{0pt} % Remove space above and below wrapfigure
% \setlength{\columnsep}{5pt} % Reduce space between text and figure
% \begin{wrapfigure}{o}{0.65\linewidth}
% \begin{figure}[htbp]
%     \centering
%      \includegraphics[width=0.65\linewidth]{figures/Cluster.pdf}
%      \caption{}
% \end{figure}
% \FloatBarrier % Forces wrapfigure to stay in place
% \setlength{\intextsep}{12pt} % Remove space above and below wrapfigure
% \setlength{\columnsep}{10pt} % Reduce space between text and figure
% The third and final task is to cluster the behavior of many model runs based on their distributions in the \textit{Annotated SOM Node Space}. 
%
% For clustering, we consider two variants.
ClimateSOM supports two types of clustering as follows.
\paragraph{Clustering model runs directly} Given the distributions of the model runs in the \textit{Annotated SOM Node Space}, we perform clustering using Earth Mover's Distance (EMD)\cite{flamary2021pot} to quantify the similarity producing a distance matrix between the distributions.
%
We then apply UMAP on the distance matrix to produce a 2D embedding and HDBSCAN \cite{mcinnes2017hdbscan} to perform the clustering.
%
While HDBSCAN can operate on the distance matrix directly, we found more consistent results using this two-step method.
%
This results in a clustering of model runs that exhibit \textit{similar} behavior and those that exhibit \textit{dissimilar} behavior, where each cluster can be represented as an aggregate distribution of the model runs in that cluster.
%
\paragraph{Clustering GCMs based on their SSP forcing} Given vector field formulation of forcing, we perform clustering of the vector field using an average cosine similarity to quantify the similarity between the vector fields.
%
Similarly, we can use UMAP and HDBSCAN to perform clustering such that we can identify GCMs that exhibit \textit{similar} behavior under the same forcing and those that exhibit \textit{dissimilar} behavior.
%
For this variant, each cluster can be represented as a vector field that represents the average behavior of the GCMs in that cluster.

A notable design implication of \ClimateSOM and its formulation is that the resulting clustering is interpretable, since it can be displayed directly in the same \textit{Annotated SOM Node Space} as distributions or vector fields.
%
% Within the broader domain of clustering algorithms and Explainable AI, there remains significant interest in the interpretability of various clustering algorithms \fix{cite}, as complex clustering workflows can often resemble a black box, making it difficult to discern the decision boundaries between clusters.
%
We consider this explainable clustering results as a key strength of \ClimateSOM.
%
% In contrast, the clustering for \ClimateSOM can be visualized directly in the \textit{Annotated SOM Node Space} as distributions or vector fields, providing a clear explanation of the clustering results.



% the SOM grid can be used as a lower-dimensional representation of the input data set.
% Once trained, the SOM  

% \subsection{Enabling interpretability of projected SOM node space}
% \subsection{Spatiotemporal data as a distribution}
% Q: Why not use DTW or other time-series distance metric?\\
% A: There are applications like some climate scenarios, where there are meaningful interpretations of the time series discarding the particular order, but considering the time series as a whole. Of course, we lose the temporal information, but we gain in simplifying the comparison.

% \subsection{Downstream tasks: Clustering and Member Comparisons}

\subsection{The Visual Interface}


To introduce the interface design for \ClimateSOM, let's consider a climate scientist with a climate ensemble dataset of monthly precipitation projections over California, expressed an anomalies from historical averages.
%
Given this ensemble, we assume the scientist in interested in the three research questions introduced in \cref{sec:problem}.

Assuming \step{S0} is complete, the trained SOM is a 2D grid of SOM nodes, where each node is a precipitation map over California.
%
Then the scientist proceeds through \step{S1} and \step{S2}.
% \subsection{\step{S1} Step 1: Anchor and \step{S2} Step 2: Annotate}

\subsubsection{SOM Node Space Editing View}
% \paragraph{Step 1: Anchor}
The main visual component for \ClimateSOM is the trained SOM and the SOM Node Space that the user iteratively edits as seen in~\cref{fig:workflow}.
%
The SOM Node Space Editing View supports the anchoring and annotating of the SOM Nodes in \step{S1} and \step{S2}.
%
First, as in~\cref{fig:workflow}-\subimage{1}, the SOM Node Space is presented to the user using the unanchored MDE result, each node depicting a precipitation map over California using a Green-Pink (high-low precipitation) colormap.
%
The interface also presents colored percentile lines in the SOM Node Space to orient the user to the mean anomaly in each SOM node.

% \begin{figure}[h]
%     \centering
%     \includegraphics[width=\linewidth]{figures/temp2.pdf}
%     \caption{Anchoring nodes reshapes the SOM Node Space in \step{S1}.}
%     \label{fig:step1}
% \end{figure}
% \label{sec:Step1and2Inteface}
% \cref{fig:Step1and2} provides an overview of the visual interface support for the two steps.
%
% First, the user is presented with \highlight{1} displaying the initial unanchored MDE result.
%
In \step{S1} the user constructs the \textit{Adjusted SOM Node Space} by anchoring and reflecting the anchors via MDE. 
%
SOM nodes in this step can be dragged around freely, and nodes that are set to be anchored are displayed with an anchor icon as seen in~\cref{fig:workflow}.
%
Furthermore, they can be inspected future via clicking on the node, displaying that node in the SOM Node Viewer (\cref{fig:Step1and2}-\subimage{1}).
%
Once they have a satisfactory \textit{Adjusted SOM Node Space} they proceed to \step{S2}.
%
This step constructs the \textit{Annotated SOM Node Space} by leveraging the LLM-empowered bidirectional query system in addition to manual visual inspection.
%
The interface supports two directions of LLM integration: forward (\cref{fig:Step1and2}-\subimage{2}) and backward (\cref{fig:Step1and2}-\subimage{3}) to support the user in defining annotations.
%
The interface provides the history of the user's interaction with the LLM in a chat-like visual UI.

Annotations as seen in~\cref{fig:Step1and2} are rendered as bounded regions in the \textit{Annotated SOM Node Space} with a user-defined label.
%
This label can be produced via the user's inspection of the bounded region in question, or with assistance from the LLM's \textasciitilde
 50-word summary of the region.
%
Once all annotations are defined and named, the \textit{Annotated SOM Node Space} is ready for analysis in \step{S3}.
%
\cref{fig:Step3} provides an overview of the visual interface support for \step{S3}.
%
\subsubsection{Model Run Inspection and Characteristic View}
% \paragraph{Model Run Inspection \subimage{1} and Characteristic View \subimage{2}}
First in pursuit of \req{R1}, the user can inspect a single model run over any given month by selecting them via a dropdown and viewing its distribution over the \textit{Annotated SOM Node Space} in the Model Run Inspection View (\cref{fig:Step3}-\subimage{1}).
%
The interface presents the 2D distributional representation of the selected model run as a KDE plot drawing three polygons as of the top 25\%, 50\%, and 75\% density regions.
%
Thus, the darkest areas of the KDE plot corresponds to the SOM Node Space regions with the highest density of realizations in the selected model run.
%
In addition, we provide the breakdown of that model run's behavior broken down by the annotations in the Characteristic View (\cref{fig:Characteristic}).
%
In the Characteristic View, the distribution for the selected model run is broken down by the annotations, presented alongside a node map glyph (\cref{fig:Characteristic}-left), ordered by its percentages in each annotation region. 
%
The node map glyph is designed to quickly direct the user to the SOM Node Space regions where the each  annotation represents.
%
Furthermore, in the case that the user is comparing two model runs as a vector field, the Characteristics View instead displays the \textit{transition} as a Sankey diagram between the annotations.


\begin{figure}[htbp]
    \centering
    \includegraphics[width=\linewidth]{Characteristic.pdf}
    \caption{The Characteristic Viewer for inspecting single model run and when comparing two model runs as a vector field.}
    \label{fig:Characteristic}
\end{figure}
\noindent
\quad\textit{Justification:}\,\, In exploring the behavior of model runs, \textbf{the primary goal for climate scientists is to understand where the \textit{hotspots} are}.
%
That is, which part of the SOM Node Space most of the realizations of the model runs are concentrated. 
%
To this end, we employ the KDE representation as the visual encoding for the distribution so that we draw the most attention to the \textit{hotspots} in the distribution. 
%
Other designs we considered included using a 3D representation and encoding the density in the height, but we found a 3D encoding less effective as it was more difficult to interpret. 
%
The Characteristic View is designed in a similar vein, to provide a quick overview of the distribution of the model run over the annotations provided in \step{S2}, and the node map glyph directing the user to the \textit{hotspots} in the distribution.

\subsubsection{Model Comparison View}
% \paragraph{Comparison Views ~\cref{fig:sidebyside} and \cref{fig:Step3}-\subimage{3}} 
To support \req{R2}, we offer two visual interfaces for comparing two model runs: via a side-by-side comparison and as a vector field.
%
First, the side-by-side comparison allows for simple visual comparison of two model runs, by placing the two distributions side-by-side in the \textit{Annotated SOM Node Space} (\cref{fig:sidebyside}).
%
By drawing the distributions over the identical SOM Node Space, the user can quickly identify the SOM Node Space regions where the densities differ.
%
However, to support the user for comparative tasks where there is a natural order to the comparison (e.g., comparing the same GCM over the historical vs. future period), we provide a vector field comparison view as well(~\cref{fig:Step3}-\subimage{3}).
\begin{figure}[htbp]
    \centering
    \includegraphics[width=0.9\linewidth]{sidebyside.pdf}
    \caption{The side-by-side comparison view for comparing two model runs.}
    \label{fig:sidebyside}
\end{figure}

\quad\textit{Justification:}\,\, 
% When considering comparison of two model runs, over discussions with our domain expert, we identified the two types of comparisons, one where is a natural ordering to the comparison and one where there is not.
% Therefore, our visual interface was designed to support both types of comparisons. 
%
Although the side-by-side comparison is simple to interpret, there were cases where the differences were subtle, or when we could infer an ordering, there was scope to provide a simpler visual encoding for the comparison.
%
We found that the vector field representation was a good fit for this task as the direction of the arrows provided a natural interpretation of the comparison as a \textit{transition}.

For the vector field formulation of the comparison, we initially used a particle advection animation similar to those used in flow visualization but found that the static vector field representation to be more simple and required less visual processing.
%
In all, the two methods for comparison were well received by our domain expert, and thus implemented in the final design.


%In \subimage{3c}, the user can inspect the breakdown of the model run over time, here for January, across the annotations provided in \step{S2}.
%%
%Secondly, to support \req{R2}, the user can compare two model runs via two visual interfaces, side-by-side \highlight{2} and via a vector field \highlight{3}.
%%
%With reference to the vector field \highlight{3} the user can understand the \textit{transition} from Selection 1 to Selection 2 with reference to the Sankey-like diagram \subimage{3e} displaying the best estimate of this transition considering the \textit{destination} of points originally in each annotation region in Selection 1.
%
\subsubsection{Clustering View}
Lastly, to support \req{R3} clustering results (\cref{ssec:clustering}) are visualized as a circuit-line diagram connecting the independent per-month clustering results(~\cref{fig:Step3}-\subimage{4}).
%
For our scientists, \textbf{in addition to understanding the results of the per-month clustering results, there is also interest in understanding the evolution of the monthly clustering results.}
%
That is, how the cluster membership of a model run changes over different months.
%
% Therefore, in addition to visualizing the clusters, we also provide the circuit-line diagram to visualize the evolution as well.
\begin{figure}[htbp]
    \centering
     \includegraphics[width=\linewidth]{step3interaction.pdf}
     \caption{
     Clicking \subimage{a} on a monthly cluster model run clustering variant \cref{fig:Step3}--\subimage{4} will select model runs represented in that cluster \subimage{b} and show its distribution over the \textit{Annotated SOM Node Space}.
    %
    On hover, only lines that pass through the hovered cluster are highlighted \subimage{c}.
    }
    \label{fig:Step3Interaction}
\end{figure}

\begin{figure*}[!t]
    \centering
    \includegraphics[width=0.95\textwidth]{step3.pdf}
    \caption{
        The visual interface to support the last step --- \step{S3} Analyze. 
        %
        Each selected model run is displayed as a distribution over the \textit{Annotated SOM Node Space} \subimage{1} and the breakdown of the distribution over the annotations are displayed in the Characteristic View \subimage{2}.
        %
        To support comparison, this main interface also supports a side-by-side comparison (\cref{fig:sidebyside}) but also as a vector field \subimage{3}.
        %
        Lastly, to support clustering of model runs or GCMs by forcing, the circuit-line diagram displays the monthly clustering results and their evolution \subimage{4}. 
        %%
        %First, in \highlight{1}, the user wishes to inspect a single model run: \req{R1}. 
        %%
        %The user can select a ensemble model run \subimage{3a} and view that model run as a distribution over the \textit{Annotated SOM Node Space} \subimage{3b}. 
        %%
        %Furthermore, given the annotations provided in \step{S2}, they are presented with a breakdown of the model run, here for January over time \subimage{3c}. 
        %%
        %The user can also inspect the list of the currently selected model \subimage{3d}. 
        %%
        %To support comparisons \req{R2}, we offer two visual interfaces, \highlight{2} and \highlight{3}, to support comparison as a side-by-side and via a vector field respectively. 
        %%
        %Of note is the comparison as a vector field in \highlight{3} where the user can inspect the comparison as a \textit{transition} from Selection 1 to Selection 2 \subimage{3e} and the interface provides a Sankey-like diagram to visualize out best estimate of transitions across the user-provided annotations. 
        %%
        %Finally \req{R3} is supported via a monthly clustering visualization \highlight{4}. 
        %%
        %\subimage{3f} and \subimage{3g} display the two variants of clusterings and the results as a circuit-line diagram connecting the per-month clustering results, placed according to 1D MDS results.
        %%
        %Interaction design and details are provided in \cref{fig:Step3Interaction} and \cref{sec:Step3Inteface}.
    }
    \vspace{-0.15in}
    \label{fig:Step3}
\end{figure*}


Each monthly cluster is represented as a box and the model runs are visualized as lines through the boxes, indicating its cluster membership for the given month.
%
The cluster boxes are colored according to the cluster's mean anomaly, and the lines according to the model run's SSP.
%
We compute a 1D MDS with intercluster distances as the distance matrix for each month, and the cluster boxes are placed according to this result.
%
Therefore, for every month, more similar clusters are placed closer together and less similar clusters further apart. 
%
In the clustering diagrams, tracing each line provides a sense of the evolution of the clustering results over the months.
%
% The lines can be filtered via the user's selection via \subimage{3h} and 
To support the investigation of the behavior of each cluster, clicking on a cluster will trigger the display of the data of that cluster as seen in \cref{fig:Step3Interaction}. 
%
Furthermore, to support the user in understanding if there are inter-month relationships between the clusters, we implement a hover interaction as well, by allowing the user to filter the lines based on those that pass through that hovered cluster.
%
A similar visual encoding supports the second variant of clustering--GCMs by forcing.
%
For this variant, each line represents GCMs instead of model runs as the GCMs are the subject of clustering.

% and we include the circuit-line diagram for this in the Appendix.



\noindent
\quad\textit{Justification:}\,\, As seen in works like \cite{jackle2016temporal} where a number of independent clustering results are the subject of analysis \textit{and} any salient trends across them, we apply a similar approach utilizing 1D MDS for placement.
%
Supporting the analysis of inter-month relationships in addition to the per-month clustering results was a key interest expressed by our domain expert, and thus we choose this design. 






\subsection{Evaluation}
% \subsection{LOCA-Downscaled CMIP6 Precipitation Ensemble}


\subsubsection{Case Studies}
\paragraph{The Data}
To demonstrate that \ClimateSOM can uncover valuable insights, we apply the workflow to a ensemble dataset of LOCA-downscaled CMIP6 precipitation projections~\cite{pierce2014statistical,pierce2015improved,eyring2016overview} over California and the Northwest United States\cite{pierce2023future}.
%
% The Localized Constructed Analogs (LOCA) statistical downscaling method \cite{pierce2014statistical,pierce2015improved } is an analog-based method used to downscale
% relatively coarse scale global climate model projections from the CMIP6 archive \cite{eyring2016overview}  to finer scale regional projections. 
%
% The LOCA North American projections were carried out at 6km spatial resolution, in which downscaling of precipitation was conducted from global model projections made by 27 CMIP6 models, 
%
These model projections were made under 3 different SSPs \cite{riahi2017shared}, from a  SSP245, a relatively low greenhouse gas emission future, SSP370, a medium high greenhouse gas emissions future, and SSP585 a high greenhouse gas emissions future. 
%
% Of the 27 models, we consider \textasciitilde 14 models for each SSP and the historical period (i.e. four time periods per GCM) for which data was available, resulting in a total of 56 ensemble members.
We consider 14 GCMs for each SSP and the historical period (i.e. four time periods per GCM) for which data was available, resulting in a total of 56 ensemble members.
%
Details of the ensemble members used in this work can be found in the Appendix.
% LOCA or Localized Constructed Analogs, is a statistical downscaling method that takes coarse resolution climate model outputs, here of CMIP6\cite{eyring2016overview}, and downscales them to a regular grid of finer resolution.

We consider two regions separately: California and the Northwestern United States, encompassing Oregon, Washington, and Idaho.
%
While daily predictions are available, we focus on monthly projections considering months from October to May, as these months contribute to \textasciitilde85\% of the annual precipitation in both spatial regions.
%
We apply two preprocessing steps. 
%
Second, we normalize the monthly anomalies by dividing them by per-month standard deviation across models, time, and space.
%
The final data, therefore, consists of monthly precipitation anomalies expressed in standard deviations from the historical average.
%
For the California case study the each ensemble model run had a of size of $780\times 151\times 165$ (time $\times$ longitude $\times$ latitude), while for the Northwestern US, it was $780\times 60\times 115$.
%
We consult \textbf{E1}, a co-author, for the California case study, and another domain expert \textbf{E2} for the Northwestern US case study.


% While this makes any analysis across different months difficult, this simplification allows us to train a single SOM across different months that streamlines the analysis.

\paragraph{LOCA-Downscaled CMIP6 over California}
\paragraph{SOM Training}
We experimented with various dimensions of the SOM grid and found that the choice had minimal impact on the results so long they were sufficiently large and selected a 30$\times$30 grid.
%
We include the impacts of the initial neighborhood radius ratio $kR$ and the neighborhood decay ratio $kS$ on the SOM in the Appendix.
%
For the California region, we selected $kR=0.03$ and $kS=0.2$ as they provided a good balance explanatory power and SOM grid smoothness.
%
This provided a SOM grid with 80.5\% of explained variance.
\paragraph{The Case Study}
We outline a case study of the precipitation ensemble over California with the \textbf{E1}, a climate scientist and co-author of this paper.
%
We will reference~\cref{fig:Dan} for this case study.
%
First, the expert goes through \step{S1}: Anchor to reorganize the SOM node space to improve the interpretability of the SOM node space.
%
They place wetter precpitation SOM nodes to the left, dryer nodes to the right, and allow the vertical axis to represent the North-South gradient of the region, placing nodes with higher precipitation in the North on the upper side of the SOM Node Space and vice versa.
%
Next, as they are familiar with the SOM nodes and their precipitation patterns, they perform \step{S2}: Annotate producing five annotations regions: (1) \textit{Statewide Wet}, (2) \textit{Wetter North - Drier South}, (3) \textit{Wetter South - Drier North}, (4) \textit{Intermediate}, and (5) \textit{Statewide Dry} (\cref{fig:Step1and2}).
%
The first question the expert wishes to answer is on the behavior of model runs that predict a wetter California in January and whether there are trends that these model run exhibit in the previous months.
% 
To this end, the expert examines the model run clustering in the circuit-line diagram and identifies two clusters of model runs in January (\req{R3}), one representing a wetter California than the other based on the color of the cluster boxes.

They then compare the two January clusters (\req{R2}), learning that the difference in the two clusters [\dist{A},\dist{B}] is primarily in the density of \textit{Statewide Wet} realizations.
%
Interestingly, the \textit{Statewide Dry} realizations are both prominently present in both clusters.
%
They visually confirm that the additional occurrences of \textit{Statewide Wet} realizations in one cluster is the driving difference between the two clusters.
%
% The expert also studies the behavior of the two clusters in December, noticing the difference here is in density in the \textit{Statewide Dry} (\req{R2}).
%
Further they learn from the circuit-line diagram that the both the drier and wetter cluster in December evolve to a wetter California in January, but realizes that it primarily evolves from a dry cluster in November considering the number of lines that pass through the respective clusters.
%
Following the dry cluster in November, they learn that this cluster is represented by heavy \textit{Statewide Dry} realizations as well as \textit{Wetter South - Drier North} realizations.
%
Thus the expert concludes that models run that predict a wetter California in January are associated with a overall drier November with a particular dryness in the North of the region, but less clear trends in December.

Another exploration that the expert wishes to undertake is of the behavior of MIROC6, a GCM, in the medium-emission scenario, SSP370, relative to the other GCMs under the same emissions scenario.
%
Exploring through different months, they observe that MIROC6 in February has a distinct pattern compared to the other GCMs in the same emissions scenario .
%
While the other GCMs over February are distributed across \textit{Statewide Wet} and \textit{Statewide Dry} realizations, MIROC6 introduces heavy \textit{Wetter North - Drier South}  and \textit{Wetter South - Drier North} realizations (\req{R1}).
%
This difference in mixture of realizations may warrant further investigation into MIROC6 - SSP370 over February to understand the reasons for the model run's distinct behavior.

\begin{figure}[!t]
    \centering
    % \includegraphics[width=\textwidth]{figures/Dan.pdf}
     \includegraphics[width=\linewidth]{Dan.pdf}
    \caption{The California case study.% with \ClimateSOM. 
     \textit{The Annotated SOM Node Space} used is shown in~\cref{fig:Step1and2}. The expert explored monthly clustering in model runs in January - here the user has hovered over \dist{A}, displaying the lines that pass through it (\req{R3}). The expert explored the cluster behavior by examining its distributions. Lastly, they also explored a single model run, \textit{MIROC6 - SSP370}, against the same SSP for all GCMs (\req{R2}).}
    \label{fig:Dan}
\end{figure}

\paragraph{LOCA-Downscaled CMIP6 over the Northwestern US}
% This section recounts a domain expert's experience using \ClimateSOM to analyze LOCA-downscaled CMIP6 precipitation ensemble over the Northwestern US.
We follow the same process to pick the SOM parameters for the Northwestern US case study, selecting $kR=0.03$ and $kS=0.2$ over a $30\times 30$ SOM grid explaining 78.8\% of the variance.
%
With consultation from an Economist, \textbf{E2}, with experience studying in climate change impacts on the Western Coast of the United States, this case study examines LOCA-downscaled CMIP6 precipitation ensemble over the Northwestern US.
%
We will reference~\cref{fig:Tom} for this case study.
%
First, the expert goes through \step{S1}: Anchor to reorganize the SOM node space to improve the interpretability of the SOM node space and \step{S2}: Annotate and visually identified regions of interest relying previous domain knowledge of precipitation patterns in the Northwestern US.
%
This yielded five annotated regions: (1) \textit{Dry Global}, (2) \textit{Wet NE}, (3) \textit{Dry SW}, (4) \textit{Wet Global}, and (5) Remainder.
% \begin{figure}[h]
%     \centering
%     \includegraphics[width=0.8\linewidth]{figures/TomSpace.png}
%     \caption{The \textit{The Annotated SOM Node Space} for the Northwestern US case study}
%     \label{fig:TomSpace}
% \end{figure}

% The expert first examines the behavior of all GCMs over the historical period in January noticing the breakdown of the selected months into the annotated regions, \textit{Dry Global} and \textit{Remainder} being the most prominent (\req{R1}).
%
% This is seen in \highlight{B}.
%
% Checking the results for the other Fall months, the expert observes the Fall vs.


First, the expert wishes to compare historical and SSP585 precipitation projections in April. 
%
Here, this transition is represented by a change from a unimodal distribution of precipitation realizations to a bimodal distribution (\highlight{C}).
%
In particular, the expert notices a shift to a wetter Northeast (of the region) under this higher emissions future, which they recognize to be a \textit{notable change} with particular significance to agriculture in the region (\req{R2}).
% \begin{figure}[h]
    % \centering
    % \includegraphics[width=\linewidth]{figures/TomBimodal.png}
    % \caption{The shift from a unimodal to bimodal distribution of precipitation realizations in April, comparing historical (left) and SSP585 (right) ensemble model runs.}
    % \label{fig:TomSpace}
% \end{figure}
%
In contrast, in November, the expert observes a movement towards a wetter Northwest (of the region) in a higher emissions future.
%
The expert therefore confirms that the effect on a higher emissions future on precipitation in the Northwestern US is drastically \textbf{different} between April and November, where April expects higher precipitation in the Northeast compared to the Northwest.
%
This is an interesting insight, as simple statistical analysis would not reveal such spatial patterns.


% with visual descriptions of the changes.
% \begin{figure}[h]
%     \centering
%     \includegraphics[width=\linewidth]{figures/TomNov.png}
%     \caption{The shift from a precipitation realizations in November, comparing historical (left) and SSP585 (right) ensemble model runs.}
%     \label{fig:TomSpace}
% \end{figure}
%
% Drilling down to a particular climate model, they observe that for ACCESS-CM2 comparing historical to a higher emission model in November, they see an increase in precpitation in the westside and coastal regions of the Northwestern US (\req{R2}).



Next, the expert wishes to explore whether there are outlier GCMs that exhibit different SSP forcing from historical to SSP585 in November. 
%
Consulting the circuit-line diagram for the GCM clustering of forcings, they observe two clusters of GCMs in November, so they click on the two clusters.
%
Here, they observe that one of the clusters (\highlight{B}), \textit{``homes in the lower right area of the SOM distribution''} thus representing GCMs that predict a drier November with drier anomalies concentrated toward the Southwest of the region.
%
On the other hand, the other cluster of GCMs (\highlight{C}), in contrast, collects toward the lower left area of the SOM distribution, and thus a wetter coastal region and wetter border with Canada, highlighting the two clusters of GCMs in its SSP forcing (\req{R3}).
%
They hypothesize that the the GCM's forcing in the first cluster may be anomalous, and thus may warrant further investigation.
%
Throughout this session, the expert was impressed with \ClimateSOM's ability to \textit{``visualize big picture patterns quickly''} and that it was \textit{``just not something (they) could do otherwise''}.




\begin{figure*}[h]
    \centering
    \includegraphics[width=\textwidth]{Tom-6.pdf}
    \caption{The Northwestern US case study with \ClimateSOM. \highlight{A}: Comparing historical and SSP585 projections in April (\req{R2}), \highlight{B},\highlight{C}: Exploring the two GCM forcing clusters in November and comparing their behavior (\req{R3}).}
    \label{fig:Tom}
\end{figure*}

\subsubsection{Evaluating the LLM integration}\label{ssec:llm_eval}
While past works demonstrate the ability of LLMs to generate SQL queries and a general text summarizer, past works on LLMs have not explicitly evaluated two spatial abilities: (1) the spatial ability to map region identifiers to a list of counties and (2) the spatial ability to summarize a list of counties to a short summary.
%
Additionally, given the possible ambiguity that may arise in how places are defined (i.e., \textit{Southern California} may mean different boundaries to different people), we conducted a short evaluation of our LLM in the above two tasks with 8 participants all familiar with the geography of California.
%
% Specifically, we designed two tasks. 

First, given a common region identifier (e.g., \textit{Southern California}), we asked the participants to evaluate how reasonable the list of counties generated by the LLM was.
%
Second, given a list of counties randomly selected from one of the textual summaries of a SOM node in the California case study, we asked the participants to evaluate how reasonable the regional summaries were.
%
Both tasks were evaluated on a 5-point Likert scale. (\textit{1: Strongly Agree, 2: Agree, 3: Neither Agree nor Disagree, 4: Disagree, 5: Strongly Disagree}).
%
We show examples of the tasks in the Appendix.

Collecting 10 of each tasks (80 responses per task), we found that for T1, the average score was 2.11 (SD=1.32) and for T2, the average score was 2.5 (SD=1.55).
%
This suggests that while the LLMs are not perfect, they are able to generate reasonable lists of counties and summaries of regions that most participants found reasonable.
%
Admittedly, in the first task there were regions where responses were more polarized, especially for vaguely defined region identifiers like \textit{Desert Region of California} and in the second task, when the list of counties were highly disconnected geographically, generating textual summaries of such regions was challenging.
%
Another limitation of the current implementation is that we use counties as the unit of analysis. 
%
However, natural language descriptions of regions do not always conform to county boundaries--some regional identifiers may span or divide multiple counties. As a result, generating finer-grained boundaries and summaries that account for these irregularities is beyond the scope of this integration.

Lastly, we also test the LLM's PostGIS generation for \ClimateSOM by test generating 20 queries for each of the three categories and manually inspecting the results.
%
We found that our LLM's generation was accurate in \textasciitilde 90\% of the cases, with the remaining \textasciitilde 10\% being cases where the LLM generated a query with incorrect syntax.
%
\textbf{In all, we find the failure cases with our LLM integration arising mainly from ambiguous region identifiers for spatial tasks, the county-level unit of analysis, and incorrect syntax generation for the PostGIS queries}.
%
This suggests a conversational interface that can allow the user to correct the LLMs when they generate unreasonable results as a future direction.

\subsection{Expert Feedback}
In addition to the two experts, \textbf{E1} and \textbf{E2}, who provided the two case studies, we also solicit feedback on \ClimateSOM with a semi-structured interview from four other domain experts, \textbf{E3} - \textbf{E6}, (total of six) presenting the case studies and the system to them.
%
These experts range in years of experience from 12 to 48 years in climate science and all have worked with climate ensemble datasets.
%
We summarize and discuss the feedback on \ClimateSOM below.

\paragraph{On the LLM integration} Since the case studies were conducted by experts who had already become familiar with the SOM nodes, thus able to rely on their past knowledge to provide annotations, we instead asked experts to test the potential of LLM integration.
%
The experts found the LLM integration particularly useful in certain circumstances. 
%
\textbf{E2} and \textbf{E3} agreed that LLM integration would be useful if the user had a specific administrative region of their interest.
%
\textbf{E1} suggested that the integration would be useful if the geospatial region in question was unfamiliar to them.
%
When visual inspection of the SOM nodes alone is insufficient to understand particular patterns in the SOM Node Space, experts \textit{agreed} that the LLM can be used to both identify regions of interest and to provide textual summaries of the nodes in those regions.
%
However, \textbf{E3} questioned the reliability of LLMs in providing sensemaking of the SOM Node Space, suggesting visual inspection of the SOM Node Space may be more reliable.
% 
Realizing a conversational interface for the LLMs to provide a way to error-correct is a future direction.

\paragraph{On the distributional nature of our data abstraction}
Throughout the case studies, we saw examples of how the distributional nature of our data abstraction of model runs in \ClimateSOM revealed important patterns.
%
Crucially, as \textbf{E1} agreed, providing a distributional view of model runs above an interpretable space via the steerable DR, annotations, and LLM-empowered sensemaking, experts can quickly identify and understand crucial characteristics of model runs.
%
% Combined with the interpretable SOM Node Space that \ClimateSOM provides via the steerable DR, annotations, and LLM-empowered sensemaking, experts can quickly identify and understand crucial characteristics of model runs. 
%
\textbf{E6} showed significant interest in the descriptive power of our representation, in particular of the vector field formulation of comparison, and noted that they had \textit{``never seen anything like it before.''}
%
Experts agreed that this distributional nature of data abstraction is a \textbf{novel way to visualize climate ensemble datasets} and a distinguishing feature of \ClimateSOM.

\paragraph{\ClimateSOM can quickly discover novel and meaningful insights} Most importantly, \textit{all experts} agreed that the insights our system uncovers in climate ensemble datasets are valuable—either because they can be identified or because identifying them manually would require substantial effort.
%
% Most importantly, \textit{all experts} agreed that the insights that our system can uncover into the climate ensemble datasets are valuable, either in that they can be identified or that it would take substantial manual effort to identify them otherwise.
%
\textbf{E2} remarked that, to generate similar insights \textit{``(their) other methods would involve generating a very large number of static plots and sifting through them.''}
% . and that they therefore \textit{`like this kind on interactive interface.`''}
%
\textbf{E1,E6} also noted the ability to investigate the monthly evolution of the clustering results in \ClimateSOM as a novel and valuable feature.
%
\textbf{E4} was impressed by \ClimateSOM and suggested that it could also be used to answer other important questions for climate models, such as whether downscaled coarse-resolution global climate models can capture the characteristics of the climate system by comparing it with high-resolution regional climate models.
%
The feedback gives us confidence that our workflow can indeed generate novel insights into climate ensemble datasets of climatological significance.


\subsection{Limitations and Future Work}
Despite the overall very positive feedback, experts also pointed out limitations.
%
First, some experts noted the \textbf{visual complexity of the workflow and interface}, though they positively reviewed the memorability of the interface (i.e., there is a learning curve to the system, but once learned, it is easy to remember how to use it).
%
Granted, the scope of tasks that \ClimateSOM attempts to address is large and complex, and thus future work will focus on improving the usability of the system with better onboarding tutorials.
%
Second, compared to common statistical or analytic methods, \ClimateSOM provides \textbf{no native metrics of statistical significance} in its findings.
%
Statistical tests are commonly used in climate science to assert anomalous patterns, for example, and the inclusion of such metrics within the \ClimateSOM framework is a future direction.
%
Third, although scalablity is not a concern for the Model Run Inspection Views, \textbf{the scalability of the Clustering View} can become a concern when considering datasets with more model runs.
%
An alternative encoding design to accommodate larger number of model runs is future work. 
%
Lastly, as introduced in~\cref{sec:problem}, the abstraction in \ClimateSOM \textbf{ignores the temporal ordering of the model runs} which is of interest to many climate scientists.

% \subsection{Discussion}
% \input{tex/8_discussion}
\subsection{Conclusion}
In this work, we introduce \ClimateSOM, a novel visual analysis workflow for climate ensemble datasets.
%
Past works show promise for visual analytics to unravel such complex datasets, but they do not address the three key tasks, exploration of a single set of model runs, comparison of two sets of model runs, and clustering of model runs and GCMs, that we identify in collaboration with a domain expert.
%
We leverage self-organizing maps, steerable DR, LLM-empowered query system to address these tasks, and present new insights into a precipitation ensemble dataset over two different spatial domains.
%
In addition, we solicit expert feedback from six scientists with experience studying such datasets and confirmed that \ClimateSOM can indeed discover climatologically significant insights that are novel or would require substantial manual effort to uncover otherwise.
%
We also discuss the limitations of \ClimateSOM and suggest future directions, in particular, the integration of statistical tests to quantify the significance of the insights generated by \ClimateSOM.
%
We believe that \ClimateSOM is a valuable tool for climate scientists to quickly and effectively explore and analyze climate ensemble datasets.
%
We encourage future work in incorporating novel VA techniques to address the unique challenges in climate science.


\section{Extending ClimateSOM: Detecting Distributional Shifts in Precipitation Co-variability}\label{sec:agu}

In this section, I continue with the \emph{ensemble variability} facet, sharpening it from visual exploration to a specific statistical question.
%
Climate-change impacts are typically assessed through mean-state changes at individual locations, yet the spatial co-variability \emph{between} regions may reorganize fundamentally under warming in ways that marginal trends alone cannot reveal.
%
Detecting such reorganization demands an explicit comparison of the distributions of spatial realizations across scenarios, something neither mean-state analysis nor EOF-based decompositions are designed to directly support.
%
To address this, I extend the ClimateSOM framework with statistical significance testing capabilities, to determine whether the \emph{distribution} of conditioned spatial precipitation patterns has shifted with statistical significance.
%
% Discovering such distributional shifts is difficult with standard analyses and here we demonstrate how ClimateSOM can be used to detect such changes under climate change.
%
I demonstrate this capability on the co-variability between California's Sacramento Bay-Delta Watershed and U.S. West Coast precipitation under LOCA2-downscaled CMIP6 projections, where I find a strengthening departure from historical relationships that I term ``decoherence,'' and the remainder of this section develops the methodology and findings in detail.
%
Although this extension remains a work in progress, I consider it a true litmus test of ClimateSOM's utility. 
%
Rather than a contrived demonstration, it applies the workflow to a climatologically substantive question, whether and how regional precipitation co-variability reorganizes under warming, and stakes its value on whether it can deliver insight that the standard analysis methods of the domain, such as mean-state trend analysis and EOF-based decomposition would struggle. 
%
It is in this confrontation with a real problem, and with the established tools climate scientists already trust, that the workflow's practical worth is genuinely put to the test.

%% Suggested section heads:
\subsection{Introduction}
\label{sec:introduction}
Climate change impacts are typically assessed through mean-state changes at individual locations or regions.
%
While such analyses provide important insights (e.g., projected wetting winters and drying springs in California~\cite{pierce2018climate,polade2017precipitation}), they overlook a critical dimension: spatial co-variability between regions may reorganize fundamentally under warming.
%
Even when marginal trends appear modest for one region, a shift in the co-variability between two regions can carry consequences that marginal trends alone do not reveal.
%
Detecting such reorganization requires an explicit comparison of distributions of spatial realizations across climate scenarios, as temporal or spatial means and correlations alone can mask compositional shifts in realized patterns.

Beyond mean-state analysis, Empirical Orthogonal Function (EOF) analysis is widely used to identify dominant modes of variability~\cite{hannachi2007empirical}, including covariability between climate modes~\cite{boschat2026covariability}. However, EOF-based approaches rely on linear combinations of orthogonal spatial patterns and may not capture nonlinear relationships or distributional shifts in realized patterns.
%
While extensions of EOFs, such as rotated EOFs~\cite{lian_evaluation_2012}, have been developed to address some of these limitations, they still primarily focus on identifying dominant modes of variability and do not directly support the distributional view that we pursue.

The value of viewing climate realizations from a distributional perspective has been shared and appreciated by past authors.
%
Pendergrass et al. examined the effects of a warmer climate considering the distribution of spatially-aggregated change in the interannual standard deviation of seasonal-mean precipitation, globally~\cite{pendergrass_precipitation_2017}.
%
Zhang et al. argues that warming produces a distributional shift in precipitation, with wider variability and more extremes~\cite{zhang_increasing_2021}.
%
Crucially, they highlight the unreliability of examining mean states alone; both because changes in mean states alone cannot inform changes in variability and because change in variability may, in fact, be more difficult to adapt to than changes in mean states.
%
A more recent work reinforced this idea by showing that day-to-day variability changes can outweigh mean-state shifts in driving precipitation extremes across many regions~\cite{nordling2025climate}.
%
However, these past works fall short of considering the distribution of the spatial realizations themselves, which we argue is critical to better understanding the impacts of climate change on regional precipitation patterns and their co-variability.

% However, we contend such analysis has been largely absent from the literature.
%
%
% framework enables rigorous hypothesis testing of whether the distribution of spatial precipitation patterns conditioned on regional climate states has shifted between historical and future periods.
To address this gap, we extend the ClimateSOM framework, building on prior applications of self-organizing maps to climate pattern analysis \cite{hewitson_self-organizing_2002, gibson_use_2017}. ClimateSOM discretizes the space of spatial climate patterns into a nonlinear, topology-preserving partition and thereby admits the distributional comparisons that EOF-based decompositions do not support. We extend this framework with a permutation-based hypothesis testing procedure utilizing Earth Mover's Distance to quantify distributional differences in spatial realizations.
% %
% Prior work has examined precipitation linkages and regional co-variability through correlation-based or mean-state analysis~\cite{}, but these approaches cannot formally test whether the composition of realized spatial patterns has shifted.
%
% Our framework enables such tests, 
% , and the stationarity assumptions underlying water resource planning.
We further develop a within-year coherence test and a clustering-based decomposition that attributes distributional shifts to compositional reorganization among identifiable pattern clusters.

We demonstrate this capability on a critical test case concerning associations between California's Sacramento Bay-Delta Watershed and U.S. West Coast precipitation.
%
The Bay-Delta is a critical hydrological hub in a persistently drought-prone region~\cite{cloern2011projected} and a focus of climate-change impact studies~\cite{rojas-aguirre_evaluating_2025}.
%
Precipitation on the U.S. West Coast has also been widely studied from its North-South gradient~\cite{dettinger1998north}, the role of atmospheric rivers in shaping seasonal totals and extremes \cite{gershunov_assessing_2017, gershunov_precipitation_2019}, and future precipitation projections in the Pacific Northwest~\cite{abatzoglou2014seasonal}.
%
However, how the co-variability between these two regions may evolve under climate change has not been systematically characterized, despite its importance for water management and climate resilience.
% \fix{ecosystem and agricultural planning across the region}.

Specifically, when the Bay-Delta experiences anomalously wet or dry water years, what spatial precipitation anomaly patterns does the broader West Coast realize, and is that relationship changing?
%
Using LOCA2-downscaled CMIP6 projections~\cite{pierce2014statistical,pierce2015improved,eyring2016overview} from 15 models under three emission scenarios~\cite{riahi2017shared} (SSP-2.45, 3.70, 5.85), we find that West Coast precipitation patterns conditioned on Bay-Delta wet and dry years are significantly departing from historical periods.
%
This departure indicates that the historical conditioning relationship between Bay-Delta and West Coast precipitation is itself changing, which we term ``decoherence''. We find that decoherence strengthens under higher emissions and over longer time horizons, most pronounced in spring months.

\subsection{Data}
\label{sec:data}
\subsubsection{LOCA2-Downscaled Climate Projections}
We analyze precipitation projections from 15 CMIP6 global climate models~\cite{eyring2016overview} (Table~\ref{tab:models}) downscaled to 6-km resolution over the contiguous United States using the Localized Constructed Analogs (LOCA2) method~\cite{pierce2014statistical,pierce2015improved,pierce2023future}.
%
LOCA2 provides monthly precipitation fields for historical (1950-2014) and future (2015-2100) periods under three Shared Socioeconomic Pathways: SSP2-4.5, SSP3-7.0, and SSP5-8.5~\cite{riahi2017shared}.
%
We also consider the first and second halves of the future period separately (2015-2057 and 2058-2100) to examine how decoherence evolves over time.
%
We consider two regions: the Sacramento Bay-Delta Watershed, a critical hydrological hub in California, and the broader U.S. West Coast region, encompassing coastal and inland areas from Washington to California (Figure~\ref{fig:regions}).

\begin{table}
\centering
\begin{tabular}{lccccc}
\hline
GCM & Variant & Historical & SSP2-4.5 & SSP3-7.0 & SSP5-8.5 \\
\hline
ACCESS-CM2       & r1i1p1f1 & \checkmark & \checkmark & \checkmark & \checkmark \\
CESM2-LENS       & r1i1p1f1 & \checkmark & \checkmark & \checkmark & \checkmark \\
CNRM-ESM2-1      & r1i1p1f2 & \checkmark & --         & \checkmark & --         \\
EC-Earth3-Veg    & r1i1p1f1 & \checkmark & \checkmark & \checkmark & \checkmark \\
EC-Earth3        & r1i1p1f1 & \checkmark & \checkmark & \checkmark & \checkmark \\
FGOALS-g3        & r1i1p1f1 & \checkmark & \checkmark & \checkmark & \checkmark \\
GFDL-ESM4        & r1i1p1f1 & \checkmark & \checkmark & \checkmark & \checkmark \\
HadGEM3-GC31-LL  & r1i1p1f3 & \checkmark & \checkmark & --         & \checkmark \\
INM-CM5-0        & r1i1p1f1 & \checkmark & \checkmark & \checkmark & \checkmark \\
IPSL-CM6A-LR     & r1i1p1f1 & \checkmark & \checkmark & \checkmark & \checkmark \\
KACE-1-0-G       & r1i1p1f1 & \checkmark & \checkmark & \checkmark & \checkmark \\
MIROC6           & r1i1p1f1 & \checkmark & \checkmark & \checkmark & \checkmark \\
MPI-ESM1-2-HR    & r1i1p1f1 & \checkmark & \checkmark & \checkmark & \checkmark \\
MRI-ESM2-0       & r1i1p1f1 & \checkmark & \checkmark & \checkmark & \checkmark \\
TaiESM1          & r1i1p1f1 & \checkmark & \checkmark & \checkmark & --         \\
\hline
\end{tabular}
\caption{Global climate models (GCMs) used in this study and their availability across SSP scenarios. Variant labels follow CMIP6 convention; absent combinations are marked with a dash and were unavailable at the time of the study.}
\label{tab:models}
\end{table}

\subsubsection{Data Preprocessing and Anomaly Calculation}
\begin{figure}
    \centering
    \includegraphics[width=0.75\textwidth]{west_coast_boundaries.pdf}
    \caption{Study region showing the Sacramento–Bay Delta Watershed (red) and the broader U.S. West Coast analysis domain (blue) used for examining precipitation co-variability.}
    \label{fig:regions}
\end{figure}

To isolate distributional changes from model biases and inter-model spread, we compute precipitation anomalies at each grid cell.
%
Here, we define the anomaly as the difference between the modeled precipitation, relative to the monthly climatology averaged across all 15 GCM models over the historical period (1950-2014).
%
This ensures our analysis detects shifts in spatial pattern distributions rather than absolute magnitude differences between models.

As in the original ClimateSOM, ensuring a pseudo-smooth set of SOM nodes that capture the full range of variability in the dataset is critical for characterizing the distribution of spatial patterns.
%
To this end, each monthly anomaly field is normalized by dividing by the standard deviation computed across all spatial locations, all SSP scenarios, all years, and all models for that specific month.
%
This normalization step is crucial for the SOM to learn meaningful patterns across months with different variability scales, preventing it from being dominated by months with higher variance.
%
Finally, we apply a signed-logarithm transformation to ensure negative and positive anomalies are equally represented in the SOM space.
%
Without transformation, the heavy positive tail of precipitation anomalies produces input vectors with extreme magnitudes that disproportionately influence SOM weight updates, causing the learned map to allocate excess capacity to representing rare extreme-wet states.
%
The signed-log compresses this asymmetry while preserving sign.

In all, for any precipitation field $p$ at month $m$ we compute the transformed anomaly $p_m'$:
\begin{linenomath*}
\begin{equation*}
    p_m' = \text{sign}(a_m) \cdot \log(|a_m| + 1), \quad a_m = \frac{p - \bar{p}_m}{\sigma_m}
\end{equation*}
\end{linenomath*}
where $\bar{p}_m$ is the historical monthly climatology (1950-2014) averaged across all 15 models for month $m$, and $\sigma_m$ is a scalar standard deviation computed across all spatial locations, SSP scenarios, years, and models for month $m$.
%
All transformed fields $p_m'$ across models, months, and time periods are pooled to train the SOM, ensuring it captures the full range of spatial patterns in the dataset.


\subsection{Methods}
\label{subsec:agu-methods}
\subsubsection{ClimateSOM Framework}
% We make use of the ClimateSOM framework with some adjustments.
% ClimateSOM~\cite{kawakami2025climatesom} was originally presented as a visual analysis tool that represents ensembles of spatial climate realizations as distributions in a two-dimensional space.
% %
% Here, we extend this framework to analyze sets of spatial realizations across different climate periods and scenarios.
%
% A self-organizing map (SOM) is trained on the transformed precipitation anomaly fields $p_m'$ (Section~\ref{sec:data}), producing a grid where each node represents a characteristic spatial pattern.
% %
% Given this SOM for any transformed precipitation field, we identify its best matching unit (BMU) in the SOM and map that node to a 2D position.
% %
% Aggregating these BMU positions across a set of realizations yields a distribution representing that set's spatial pattern composition.
% %
% Formally, for a set with $n$ transformed spatial fields $F=[f_1, f_2, \ldots, f_n]$, we obtain:
% \begin{equation}
% S=\{s_1,s_2,\dots,s_n\}, \,\,\, s_i=\text{proj}(\text{BMU}(f_i))
% \label{eq:climatesom_projection}
% \end{equation}
% where $\text{BMU}(f_i)$ is the best matching SOM node for transformed field $f_i$, and $\text{proj}(\cdot)$ maps the node to 2D coordinates.
% %
% In this work, we use a $30\times 30$ SOM grid, following the configuration established in the original ClimateSOM work~\cite{kawakami2025climatesom}. 
% \subsubsection{SOM Projection}
As identified, the original 2D SOM grid often contains distortions that we wish to mitigate.
%
The $\text{proj}(\cdot)$ function maps each SOM node to 2D coordinates that serve as the ground space for our EMD computation.
%
In this work, we use UMAP~\cite{mcinnes2018umap}, departing from the original ClimateSOM implementation, which used Minimum-Distortion Embedding (MDE)~\cite{agrawal2021minimum}.
%
MDE is a general framework that generalizes UMAP and t-SNE as special cases, and the additional anchoring constraints it exposes are not needed here, so we use UMAP as a simpler instance of the same family.

Local fidelity is the appropriate property for our application because EMD aggregates the cost of moving mass over short hops, and the climate-change signal we test takes the form of many such local redistributions in pattern composition rather than wholesale traversal across the pattern space.
%
Global distortions introduced by UMAP are therefore of limited concern, as long-range transports are less prevalent in our dataset. 
%
Instead, individual displacements are local, and their accumulation across realizations is what we report as decoherence.
%
The UMAP projection is computed once on the full SOM codebook, inspected for validity, and held fixed across all comparisons and permutations, so the test statistic depends only on group assignments and not on embedding stochasticity.

\subsubsection{EMD-Based Hypothesis Testing}
We extend ClimateSOM with a permutation test using Earth Mover's Distance (EMD) to formally compare distributions of spatial climate patterns between different sets of realizations.
% \shortnote{What if there's a different number?}
%
Recall that the SOM is trained once on the full dataset (all historical and future fields, all models and scenarios), ensuring a consistent pattern space for all comparisons.
%
Given two sets of spatial realizations with BMU distributions $S_1 = \{s_1^{(1)}, s_2^{(1)}, \ldots, s_{n_1}^{(1)}\}$ and $S_2 = \{s_1^{(2)}, s_2^{(2)}, \ldots, s_{n_2}^{(2)}\}$ in the UMAP-reduced ClimateSOM space, we compute the observed EMD:
\begin{equation}
d_{\text{obs}} = \text{EMD}(S_1, S_2)
\end{equation}
computed exactly via the network simplex solver in POT~\cite{flamary2021pot}. 
%
Here we use Euclidean distance between coordinates.
%
Additionally, we note that we solve the balanced optimal transport problem, which assumes equal mass in the compared distributions.
%
To satisfy this assumption, we equally distribute mass across all points in any set $S_i$ such that for any $S_i$ the total mass sums to 1.
%
This allows us to test distributions of arbitrary size under our framework.
%
The UMAP projection is computed once on the full SOM pattern space and held fixed across all permutations, ensuring the test statistic depends only on group assignments. 
% As an additional validation, we compare our EMD-based results against a standard two-sample t-test on the BMU coordinates and found qualitatively consistent significance patterns.

To assess statistical significance, we test the null hypothesis $H_0$ -- that the two distributions are drawn from the same underlying pattern distribution.
%
We perform a permutation test by pooling $S_1 \cup S_2$, randomly shuffling group labels, and recomputing EMD for each permutation.
%
Let $d_{\text{perm}}^{(k)}$ denote the EMD for the $k$-th permutation.
%
The p-value is calculated as:
\begin{equation}
p = \frac{1}{K} \sum_{k=1}^{K} I(d_{\text{perm}}^{(k)} \geq d_{\text{obs}})
\end{equation}
where $I(\cdot)$ is the indicator function.
%
We empirically find that $K=5000$ provides stable estimates and employ that here.
%
Lastly, we use $\alpha = 0.05$ as the significance threshold.
%
The validity of this permutation test rests on \emph{exchangeability} under $H_0$: relabeling realizations must be equally likely when the two sets are drawn from the same pattern distribution. We obtain this by treating each (model, water year) realization as an independent draw from the population of plausible climate states under its scenario. We acknowledge this independence assumption as a limitation, since realizations are correlated both across models, which share components and physics, and across neighboring water years; the resulting reduction in effective sample size may render our p-values mildly anti-conservative.

This approach detects both compositional reorganization (changes in which spatial patterns are realized) and pattern intensification, with one caveat.
%
The signed-log transformation compresses magnitude, so BMU assignment is pattern-dominated; intensified patterns shift to a more extreme neighboring node when one exists.
%
At the wettest and driest SOM extremes, however, no more-extreme nodes exist, so intensification of those patterns remains invisible to the test.
%
Significant decoherence therefore indicates a shift in the spatial pattern composition associated with a given regional regime, while non-significance does not strictly rule out intensification confined to the SOM extremes.

\subsubsection{Application to West Coast Precipitation Co-variability}

We apply this framework to test whether the relationship between Bay-Delta precipitation regime and West Coast spatial precipitation patterns may reorganize under climate change.
%
Specifically, we examine whether the spatial patterns realized across the broader West Coast during anomalously wet or dry Bay-Delta water years differ systematically between historical and future periods.

We classify each water year as wet or dry based on Bay-Delta accumulated precipitation, applying the same historical thresholds (defined below) to both historical and future periods so that wet/dry labels are consistent across all comparisons.
% %
% For each water year, we calculate the total accumulated precipitation in the Bay-Delta region.
%
Wet years are defined as those exceeding the 75th percentile of the historical distribution (1950-2014, pooled across all 15 models), and dry years fall below the 25th percentile.
%
For months within wet water years, we extract the corresponding West Coast spatial precipitation patterns from the historical period and compare them to West Coast patterns from months within wet water years in the future period.
%
All 15 climate models are pooled together in each comparison.
%
We repeat this procedure separately for dry water years.
%
Comparisons are conducted across October through May, three SSP scenarios (SSP2-4.5, SSP3-7.0, SSP5-8.5), and two future time periods (early: 2015-2057, late: 2058-2100).
%
% A significant p-value indicates that the West Coast precipitation response to a given Bay-Delta regime has fundamentally changed between historical and future climates.
% \subsection{Comparison and verification with standard analysis approaches}
% To assess the validity and robustness of our EMD-based hypothesis testing result, we compare our results

\subsection{Results}
\label{sec:results}
\begin{figure}
    \centering
    \includegraphics[width=0.9\textwidth]{som_grid_7x7.pdf}
    \caption{An evenly selected 7$\times$7 grid of SOM nodes from the trained ClimateSOM, showing the spatial precipitation anomaly patterns associated with each node. Values are in the log-standardized anomaly space used for training as  described in Section~\ref{subsec:agu-methods}. Broadly the SOM extracts two axes of change: (1) the North-South axes of wetness and dryness and (2) overall wetness and dryness across the entire West Coast. Each SOM node is projected via
    UMAP~\cite{mcinnes2018umap} and its projected position is used to represent any precipitation anomaly pattern in the dataset.
    \label{fig:som_grid}}
\end{figure}
% \subsection{Changes in Bay-Delta Wet and Dry Year Frequency}
%
% Before examining distributional changes in spatial patterns, we first test whether the frequency of Bay-Delta wet and dry water years is itself changing under climate scenarios.
% %
% Table~\ref{tab:wy_frequencies} shows the proportion of water years classified as dry (below 25th percentile) or wet (above 75th percentile) based on Bay-Delta accumulated precipitation across all 13 models for each scenario and time period.
% %
% Fisher's exact test reveals significant shifts in these proportions under higher emissions and in the late-century period.
% %
% SSP2-4.5 and SSP3-7.0 in the second half of the century show both a decrease in dry year frequency and a significant increase in wet year frequency, with wet years comprising 35.89\% and 32.75\% of water years respectively compared to the historical 25.00\%.
% %
% However, most notably, such changes do not occur under the highest emissions scenario (SSP5-8.5), where dry and wet year frequencies remain statistically indistinguishable from historical proportions in both future periods.
% %
% \shortnote{Whats this about?}
%
%
% \begin{table}
% \centering
% \caption{Frequency of Bay-Delta dry and wet water years across scenarios and time periods. Bold values differ significantly from historical proportions (Fisher's exact test, $p < 0.05$).}
% \label{tab:wy_frequencies}
% \begin{tabular}{c c c c}
% \hline
% Scenario & Period & Dry Model-Years (\%) & Wet Model-Years (\%) \\
% \hline
% Historical & 1950--2014 & 25.10 & 25.00 \\
% \hline
% SSP2-4.5 & Early (2015--2057) & 21.59 & 26.79 \\
%          & Late (2058--2100) & \textbf{20.21} & \textbf{35.89} \\
% \hline
% SSP3-7.0 & Early (2015--2057) & 25.16 & 26.62 \\
%          & Late (2058--2100) & \textbf{23.52} & \textbf{32.75} \\
% \hline
% SSP5-8.5 & Early (2015--2057) & 24.13 & 27.10 \\
%          & Late (2058--2100) & 24.58 & 29.46 \\
% \hline
% \multicolumn{4}{l}{\footnotesize Dry: $<$25th percentile; Wet: $\geq$75th percentile of historical distribution}
% \end{tabular}
% \end{table}

\subsubsection{Historical Model Evaluation Against Observations}

Before examining future changes, we evaluate whether the CMIP6 models capture observed historical co-variability patterns by comparing each model's historical period (1950-2014) against Livneh gridded observational data (1950-2014)~\cite{pierce2021extreme}, both as an (1) validation of each GCM to adequately \emph{replicate} historical patterns and (2) as a validation of our methodology.
%
We apply the same EMD-based permutation test to compare West Coast precipitation pattern distributions conditioned on Bay-Delta wet and dry years between Livneh observations and each GCM's historical simulation.
%
Figure~\ref{fig:livneh_unsplit_pvalues} shows p-values for these comparisons across months and conditioning states.
%
Most models show no significant decoherence from observations for most months, indicating that the models generally capture the historical co-variability structure between Bay-Delta regime and West Coast spatial patterns.
%
This provides confidence that detected future changes represent genuine reorganization signals rather than model artifacts.
%
% \begin{figure}
%     \centering
%     \includegraphics[width=\textwidth]{livneh_pvalue_table.pdf}
%     \caption{P-values from EMD permutation tests comparing Livneh observed precipitation patterns (1950-2011) against each CMIP6 model's historical simulation (1950-2014). Non-significant p-values indicate models that show no significant signs of deviations.}
%     \label{fig:livneh_pvalues}
% \end{figure}

\begin{figure}
    \centering
    \includegraphics[width=\textwidth]{livneh_unsplit_pvalue_table.pdf}
    \caption{P-values from EMD permutation tests comparing Livneh unsplit observed precipitation patterns (1950-2014) against each CMIP6 model's historical simulation (1950-2014). Non-significant p-values indicate models that show no significant signs of deviations.}
    \label{fig:livneh_unsplit_pvalues}
\end{figure}

\subsubsection{Welch t-test based hypothesis testing results}\label{ssec:t_test_results}
To assess the validity of our EMD-based decoherence test, we compare its results against a standard parametric approach using Welch's t-test on mean West Coast precipitation anomalies.
%
Welch's variant does not assume equal variances, an appropriate choice given the differing sample sizes and variance characteristics across periods and scenarios.

The two tests address related but distinct questions.
%
The t-test detects shifts in the mean magnitude of anomalies between periods, whereas the EMD test detects shifts in the distribution of spatial patterns themselves.
%
Agreement between the two provides converging evidence of distributional change; disagreement would highlight cases where pattern composition shifts without a corresponding shift in mean intensity.

We compute the t-test for each combination of month, scenario, and conditioning state, comparing all grid-point anomalies in the West Coast domain between historical and future periods.
%
Results appear alongside the EMD results in Figure~\ref{fig:combined_pvalues}.
%
Broadly, the two tests show consistent significance patterns, supporting the validity of our EMD-based approach.
%
Across all 96 test cells (3 SSP scenarios $\times$ 2 future half-periods $\times$ 2 Bay-Delta conditioning states $\times$ 8 months), the two tests reach the same significance verdict at $p < 0.05$ in 71 cases (74\%), as summarized in Table~\ref{tab:emd_welch_agreement}.
%
Critically, the disagreements are highly asymmetric: in 22 of the 25 discordant cells the EMD test detects a significant distributional shift where the Welch t-test does not, whereas the reverse occurs in only 3 cells.
%
This asymmetry is consistent with our central premise, that the spatial composition of precipitation patterns can reorganize without a detectable change in mean intensity, and that such reorganization is visible to a distributional test yet invisible to a mean-based one.
%
\begin{table}[t]
    \centering
    \begin{tabular}{lccc}
        \toprule
        & \multicolumn{2}{c}{Welch t-test} & \\
        \cmidrule(lr){2-3}
        & Significant & Not significant & Total \\
        \midrule
        EMD significant     & 19 & 22 & 41 \\
        EMD not significant & 3  & 52 & 55 \\
        \midrule
        Total               & 22 & 74 & 96 \\
        \bottomrule
    \end{tabular}
    \caption{Agreement between the EMD-based permutation test and the Welch t-test across all 96 test cells (3 SSP scenarios $\times$ 2 future half-periods $\times$ 2 Bay-Delta conditioning states $\times$ 8 months), with each cell classified as significant or not at $p < 0.05$. The tests agree in 71 of 96 cells (74\%): 19 jointly significant and 52 jointly non-significant. Of the 25 disagreements, 22 are cells where the EMD test detects a distributional shift that the mean-based Welch test misses, and only 3 are the reverse.}
    \label{tab:emd_welch_agreement}
\end{table}
%
The following sections focus on the EMD-based analysis.
% \begin{figure}
%     \centering
%     \includegraphics[width=\textwidth]{welch_pvalue_table.pdf}
%     \caption{P-values from Welch t-tests comparing mean West Coast precipitation anomalies conditioned on Bay-Delta wet and dry years between historical (1950-2014) and future periods (2015-2057, 2058-2100) across three SSP scenarios }
%     \label{fig:t_test_pvalues}
% \end{figure}

\subsubsection{EMD-based hypothesis testing results}
In the following section, we evaluate the future vs historical West Coast precipitation distributional differences using our ClimateSOM extended EMD-based method, instead of considering the mean of all precipitation anomaly patterns as was performed above ~\ref{ssec:t_test_results}.
% 
Again the precipitation patterns are conditioned upon Bay-Delta watershed precipitation being wet or dry.
%
Our permutation tests reveal significant decoherence in West Coast precipitation pattern distributions conditioned on Bay-Delta regime from historical to future periods (Figure~\ref{fig:combined_pvalues}).
%
The magnitude and timing of decoherence depend systematically on emission scenario, time horizon, conditioning state (wet vs. dry), and month.

% \begin{figure}
%     \centering
%     \includegraphics[width=\textwidth]{pvalue_table.pdf}
%     \caption{P-values from Earth Mover's Distance permutation tests (1,000 permutations) comparing distributions of West Coast spatial precipitation patterns conditioned on Bay-Delta wet and dry years between historical (1950-2014) and future periods. Results span 13 CMIP6 models aggregated, three SSP scenarios (rows), two future time periods (early: 2015-2057, late: 2058-2100), Bay-Delta wet and dry conditioning states, and eight months (Oct-May). Significant decoherence ($p < 0.05$) strengthens with higher emissions, later time horizons, and varies systematically by month and conditioning state.}
%     \label{fig:pvalues}
% \end{figure}

\begin{sidewaysfigure}
    \centering
    \includegraphics[width=\textwidth]{combined_pvalue_table.pdf}
    \caption{Combined p-values from Welch t-tests and EMD-based permutation tests comparing historical and future distributions of West Coast precipitation patterns conditioned on Bay-Delta wet and dry years. }
    \label{fig:combined_pvalues}
\end{sidewaysfigure}

Decoherence strengthens from early-century to late-century periods.
%
In the early future period (2015-2057), Bay-Delta dry years show emerging decoherence primarily in October and April-May under higher emissions (SSP5-8.5: $p = 0.000$ in October, $p = 0.017$ in April), while wet years show weaker signals with only scattered significant months.
%
Table~\ref{tab:emd_sig_counts} summarizes this progression by tallying, for each scenario and conditioning regime, the number of months (out of eight, October--May) that exhibit significant distributional divergence ($p < 0.05$) under the EMD test.
%
The tally of significant months more than doubles from the early to the late period, from 13 to 28 summed across all scenarios and conditioning states, confirming that decoherence broadens across the wet season as the century progresses rather than remaining confined to a handful of months.
%
The signal also grows monotonically with emissions within nearly every period-regime block and is consistently stronger in dry-year than wet-year conditioning during the first half, with wet years catching up only in the late century.

\begin{table}[t]
    \centering
    \begin{tabular}{lcccc}
        \toprule
        & \multicolumn{2}{c}{First-Half (2015--2057)} & \multicolumn{2}{c}{Second-Half (2058--2100)} \\
        \cmidrule(lr){2-3} \cmidrule(lr){4-5}
        Bay-Delta Conditioning & Dry & Wet & Dry & Wet \\
        \midrule
        SSP2-4.5 & 3 & 1 & 5  & 4  \\
        SSP3-7.0 & 2 & 1 & 4  & 4  \\
        SSP5-8.5 & 4 & 2 & 6  & 5  \\
        \bottomrule
    \end{tabular}
    \caption{Number of months (out of eight, October--May) with significant EMD-based distributional divergence ($p < 0.05$) between historical and future West Coast precipitation pattern distributions, by SSP scenario, future half-period, and Bay-Delta conditioning regime.}
    \label{tab:emd_sig_counts}
\end{table}

%
The late-century period (2058-2100) reveals widespread decoherence across both conditioning states.
%
For Bay-Delta dry years, fall and early winter months (October-January) show strong significance across all SSP scenarios, particularly under SSP3-7.0 and SSP5-8.5, and spring months (April-May) maintain significance across scenarios.

Bay-Delta wet years in the late period exhibit the strongest reorganization signals.
%
Under SSP5-8.5, late winter through spring (March-May) shows highly significant decoherence ($p = 0.000$ for all three months), with additional shifts in November ($p = 0.003$) and February ($p = 0.019$).
%
SSP3-7.0 shows similar spring decoherence (April-May: $p = 0.000$) and significant shifts in November ($p = 0.007$) and February ($p = 0.000$).
%
Additionally, SSP2-4.5 shows significant late-century wet-year decoherence in March-May.

Winter core months (December-January) generally show weaker signals, particularly for wet years.
%
This may indicate that mid-winter co-variability is undergoing intensification of existing spatial patterns rather than compositional reorganization; such changes would not be detected by our BMU distribution comparison, as intensified patterns would still map to the same SOM nodes.
%
The asymmetry between wet and dry year responses, especially pronounced in the late century, indicates that the Bay-Delta relationship with West Coast patterns is reorganizing in regime-dependent ways rather than through simple uniform shifts.
%
These results demonstrate that historical co-variability structures are fundamentally breaking down under climate change, with the rate and character of reorganization varying systematically by emission pathway, time horizon, and regional conditioning state.

% \begin{figure}
%     \centering
%     \includegraphics[width=\textwidth]{median_pvalue.png}
%     \caption{Median p-values across MIROC6, EC-Earth3-Veg, KACE-1-0-G, ACCESS-CM2, IPSL-CM6A-LR, INM-CM5-0 comparing pooled historical and future ssp370}
%
% \end{figure}


% is not enough to detect significant differences. I think the best way to show this is to just show the distribution of p-values across models for each cell in Figure 4, but even then I think it will be hard to interpret. Let me know if you want me to try that or if you have other ideas.}

\subsubsection{Visualizing the Overall Spatial Character of Reorganization}\label{ssec:cumulative_change}

Given the significant decoherence detected in the previous section, we now visualize \textit{what} is changing spatially.
%
To this end, we compute cumulative change fields where each SOM node's spatial pattern is weighted by the shift in its occupancy frequency.
%
For each comparison between distributions $S_1$ (future) and $S_2$ (historical), we calculate the proportion of samples matching each SOM node $i$ in both distributions. 
\begin{align*}
    p_{1,i} &= \text{count}(S_1 \text{ matches to node } i)/|S_1| \\
    p_{2,i} &= \text{count}(S_2 \text{ matches to node } i)/|S_2|
\end{align*}
The cumulative change field is then
\begin{linenomath*}
\begin{equation*}
    \text{Cumulative Change} = \sum_{i} (p_{1,i} - p_{2,i}) \cdot W_i
\end{equation*}
\end{linenomath*}
where $W_i$ is the spatial precipitation pattern associated with SOM node $i$.
%
Positive values indicate locations whose expected precipitation under the given conditioning is higher in the future than historically; negative values, lower.
%
We note that the field is a relative change between historical and future distributions and does not characterize the absolute wetness or dryness of locations as projected in the CMIP6 catalog.


\begin{sidewaysfigure}
    \centering
    \begin{subfigure}{0.49\textwidth}
        \centering
        \includegraphics[width=\textwidth]{change_mosaic_grid_first_dry.pdf}
        \caption{Dry Bay-Delta years, early future (2015--2057)}
        \label{fig:change_mosaic_grid_first_dry}
    \end{subfigure}
    \hfill
    \begin{subfigure}{0.49\textwidth}
        \centering
        \includegraphics[width=\textwidth]{change_mosaic_grid_second_dry.pdf}
        \caption{Dry Bay-Delta years, late future (2058--2100)}
        \label{fig:change_mosaic_grid_second_dry}
    \end{subfigure}

    \vspace{0.75em}

    \begin{subfigure}{0.49\textwidth}
        \centering
        \includegraphics[width=\textwidth]{change_mosaic_grid_first_wet.pdf}
        \caption{Wet Bay-Delta years, early future (2015--2057)}
        \label{fig:change_mosaic_grid_first_wet}
    \end{subfigure}
    \hfill
    \begin{subfigure}{0.49\textwidth}
        \centering
        \includegraphics[width=\textwidth]{change_mosaic_grid_second_wet.pdf}
        \caption{Wet Bay-Delta years, late future (2058--2100)}
        \label{fig:change_mosaic_grid_second_wet}
    \end{subfigure}
    \caption{Cumulative change fields of West Coast precipitation relative to the historical period (1950--2014), aggregated across 15 CMIP6 models by month (columns) and SSP scenario (rows within each panel), conditioned on Bay-Delta dry years (top) and wet years (bottom) for the early (left) and late (right) halves of the future period. Positive (teal) values indicate locations wetter in the future than historically under the given conditioning, and negative (brown) values drier; bordered cells denote months with significant distributional divergence ($p < 0.05$).}
    \label{fig:change_mosaic_grid}
\end{sidewaysfigure}

%
%
%
%

\subsubsection{Shifting seasonal precipitation in the U.S. West Coast, conditioned on Bay-Delta wet and dry years}
\paragraph*{Bay-Delta Dry Years}
Bay-Delta dry years show the strongest decoherence in fall (October-November) and spring (April-May), particularly under higher emissions and in the late future period.
%
The cumulative change fields reveal that these changes are primarily driven by a shift of precipitation away from the Bay-Delta Watershed toward the Pacific Northwest, with the exception of May in both periods where Bay-Delta dryness co-occurs with an overall dry U.S. West Coast (Figures~\ref{fig:change_mosaic_grid_first_dry} and~\ref{fig:change_mosaic_grid_second_dry}).

\paragraph*{Bay-Delta Wet Years}
Bay-Delta wet years show modest decoherence in the early future period, but exhibit more systematic and stronger decoherence in the late future period, particularly under SSP5-8.5.
%
The most notable changes occur in spring months (MAM), especially May, where Bay-Delta wet water years associate with a shift toward drier conditions in California and the Sierra Nevada, and wetter conditions in the Pacific Northwest (Figure~\ref{fig:change_mosaic_grid_second_wet}).
%
% This suggests that even during anomalously wet water years in the Bay-Delta, spring months may experience relative dryness compensated by wetter conditions in other months of the water year \fix{this needs to be verified}.
%
Similar trends appear in November under SSP3-7.0 and SSP5-8.5 in the late future period.


\subsubsection{Within-year coherence of cumulative changes}
The cumulative change fields shown in Figures~\ref{fig:change_mosaic_grid_first_dry},~\ref{fig:change_mosaic_grid_second_dry},~\ref{fig:change_mosaic_grid_first_wet}, and~\ref{fig:change_mosaic_grid_second_wet} represent aggregate patterns computed across all years within each period-condition.
%
However, a critical question remains unanswered. 
%
Do these spatial patterns actually manifest together within individual water years, or do they emerge only when averaging across the entire period?
%
To address this, we test whether the cumulative change patterns are coherently expressed in specific water years at rates higher than expected by chance.

For each (month $m$, condition $c$), we first compute the historical conditional mean anomaly field $\bar{H}_{m,c}$ by averaging historical anomalies for that month and condition.
%
For each future month-model-year under condition $c$, we then compute the delta from this baseline:
\begin{equation*}
\delta_{m,y,k,c} = a_{m,y,k} - \bar{H}_{m,c}
\end{equation*}
This delta is directly comparable to the cumulative change field, since both are deltas relative to historical conditional behavior.
%
The subtraction also removes confounding from the historical conditional pattern (for example, the typical wet-year pattern that already existed historically) bleeding into the comparison.
%
We then calculate the Pearson correlation between each per-month delta $\delta_{m,y,k,c}$ and the corresponding cumulative change field, and average these monthly correlations across Oct-May for each model-year, yielding a single coherence metric per model-year.
%
A model-year with high mean correlation indicates that multiple months within that year exhibit shifts spatially aligned with the cumulative change field.

We use 0.3 as a threshold for identifying model-years with meaningful coherence.
%
While modest, this threshold provides a conservative indicator of systematic pattern expression.
%
We count the number of model-years exceeding this threshold within each period-condition as our test statistic.

To determine whether the observed counts are statistically significant, we perform another permutation test.
%
For each month within a period-condition, we randomly reassign all correlation values for that month across different year-model combinations.
%
This shuffling preserves the distribution of correlation values for each month but removes any temporal coherence within individual water years.
%
We recalculate model-year mean correlations from these shuffled values and count how many exceed 0.3, repeating this process to build a null distribution.
%
Under the null hypothesis that the cumulative changes are purely an aggregate artifact with no systematic expression in individual years, the observed count would be consistent with the null distribution.
%
Observed counts significantly exceeding the null distribution indicate that specific water years do indeed exhibit the cumulative change patterns coherently across multiple months.
% we perform a permutation test on the count of year-model groups with mean correlation $> 0.3$ between the change fields for each period-condition.
\begin{table}
    \centering
\begin{tabular}{l r r r r}
\hline
Period-Condition & Observed Count ($>0.3$)  & Null Mean & Null SD & p-value \\
\hline
First - Dry & 26  & 19.83 & 3.87 & 0.0778 \\
First - Wet & 51  & 46.48 & 5.15 & 0.2190 \\
Second - Dry & 93  & 76.12 & 5.62 & \textbf{0.0018} \\
Second - Wet & 113 & 91.91 & 6.45 & \textbf{0.0012} \\
\hline
\end{tabular}
\caption{Results of permutation test for co-occurrence of cumulative change fields with monthly anomaly fields across model-years for each period-condition ($\alpha=0.05$). Significant p-values indicate that the observed count of model-years with mean correlation $>0.3$ is greater than expected under the null distribution generated by randomly reassigning correlation values within each month across year-model combinations.}
\label{tab:per_period_pvals}
\end{table}
%
We show the results in Table~\ref{tab:per_period_pvals}.
%
The results demonstrate that the co-occurrence of cumulative changes is not merely an artifact of aggregating across entire periods, but rather reflects patterns coherently expressed in individual model-years at rates significantly higher than expected by chance.
%
This is particularly evident in the second half of the century for both dry and wet Bay-Delta years, where observed counts substantially exceed null expectations.
%
% Figure~\ref{fig:annual_cooccurrence} shows the number of model-years with mean correlation $>0.3$ for each year within each period-condition, illustrating the temporal evolution of within-year coherence.

% \begin{figure}
%     \centering
%     \includegraphics[width=0.7\textwidth]{annual_cooccurence.png}
%     \caption{Annual count of model-years with mean correlation $>0.3$ between monthly anomaly fields and cumulative change fields for each period-condition. The increase in counts in the second half of the century visually supports the significant co-occurrence results from the permutation test.}
%     \label{fig:annual_cooccurrence}
% \end{figure}

% \subsection{Discussion}
% \label{sec:discussion}
% \shortnote{To add: Implications, Limitations}

\subsection{Continued Work}
\label{subsec:agu-conclusion}

This section forms the last of three arms in our extension of ClimateSOM.
%
First, we extended the workflow with a statistical framework, an EMD-based permutation test that determines whether the distribution of conditioned spatial precipitation patterns has shifted under a changing climate.
%
Second, we established the validity of this framework by validating it against a traditional approach, Welch's t-test on mean precipitation, and found the two in broad agreement.
%
Third, in this final arm, we develop how to visualize the change that the test detects; this work is ongoing, and we describe it here to outline where we are currently taking the extension.

Our aim is to visualize the spatial reorganization that our EMD test has detected.
%
The cumulative change fields above give a first, aggregate view of this reorganization, and we are now developing methods to resolve that view further, decomposing the aggregate change into the specific spatial pattern transitions that compose it, so that the detected reorganization can be visualized as a set of interpretable shifts between precipitation patterns rather than a single blended field.

\subsubsection{Decomposing the Spatial Reorganization}\label{ssec:decomposing_reorganization}
The cumulative change field $\sum_i (p_{1,i} - p_{2,i}) W_i$ defined in Section~\ref{ssec:cumulative_change} aggregates contributions from many SOM nodes and therefore does not isolate the specific pattern transitions responsible for the observed reorganization.
%
To resolve this, we are developing a clustering-based decomposition (Algorithm~\ref{alg:decomposition}) that couples optimal transport with density-based clustering of SOM-projected samples in the ClimateSOM space (Figure~\ref{fig:agu_som_space}), quantifying mass redistribution between historical and future periods at the level of pattern clusters.
%
We select HDBSCAN over $k$-means and hierarchical clustering because it does not require the number of clusters to be specified \textit{a priori} and accommodates clusters of varying shapes and densities. Points classified by HDBSCAN as noise are excluded from the decomposition.

\begin{figure}[t]
    \centering
    \includegraphics[width=\textwidth]{AGU_SOM_Space.pdf}
    \caption{The two-dimensional ClimateSOM space over which the decomposition operates. Each point is a SOM node placed at its projected position, and the callout shows the log-standardized precipitation anomaly pattern associated with one node. Every sample is located in this space by its best-matching unit, so a set of samples becomes a density over the space that the decomposition then clusters and transports.}
    \label{fig:agu_som_space}
\end{figure}

The reduced transport plan $\tilde{P}$ quantifies the mass redistributed between each pair of pattern clusters; however, redistributed mass alone is an insufficient measure of the magnitude of spatial change.
%
A large transfer between two similar clusters constitutes a minor reorganization, whereas a smaller transfer between dissimilar clusters may contribute substantially to the cumulative change field.
%
To rank transitions by their contribution to the overall reorganization, we define an \textit{importance score} for each cluster pair as the product of (1) the redistributed mass $\tilde{P}_{ab}$ and (2) the mean absolute error (MAE) between the cluster centroids in the SOM prototype space.
%
The MAE between centroids in the SOM prototype space directly measures the dissimilarity between the underlying spatial precipitation patterns. High importance scores then identify transitions in which substantial mass is exchanged between markedly distinct patterns.

Although unnormalized importance scores have different totals across comparisons, reflecting different magnitudes of total reorganization, expressing each as a percentage of its comparison's total yields a compositional measure: the fraction of that comparison's spatial reorganization attributable to a given cluster transition.
%
This normalized score lets us rank the cluster transitions driving the observed changes and compare their contributions across comparisons, even when the underlying comparisons differ in scale or sample size. In practice, we examine the top-ranked cluster transitions per comparison to identify the dominant pattern shifts.
%
Figure~\ref{fig:agu_decomposition} illustrates the decomposition for an example comparison: historical versus SSP5-8.5 March precipitation in the late-century period (2058--2100), conditioned on Bay-Delta wet years.
%
We select this comparison because the EMD test flags it as a significant distributional shift (Figure~\ref{fig:combined_pvalues}, $p = 0.001$), so a genuine reorganization is known to be present for the decomposition to resolve.
%
The before and after distributions are each partitioned into clusters, and the optimal transport plan redistributes probability mass between them to isolate the dominant cluster transitions.
%
Figure~\ref{fig:agu_centroid_weights} then renders, for the five highest-scoring transitions, the source and target cluster centroids as spatial precipitation patterns, making the specific reorganization each transition represents directly legible.

Conditioned on Bay-Delta wet March years in the late century under SSP5-8.5, the five highest-scoring transitions share a consistent theme: a loss of precipitation in the south giving way to gains in the north.
%
The top transition (Flow~1, importance score $17.35\%$) relaxes an extreme-dry Sierra Nevada and Northern California source pattern toward a more moderate dryness.
%
Flow~2 ($14.49\%$) loses a wet Southern California pattern to a slightly dry Southern California and Sierra state.
%
Flow~3 ($12.49\%$) introduces a new wet Northern California and Sierra pattern.
%
No single transition dominates, but instead the reorganization spreads across several comparably-weighted flows that together shift wet-year March precipitation away from Southern California and toward Northern California and the Pacific Northwest.
%
\begin{algorithm}[h!]
\caption{Clustering-based decomposition of distributional reorganization}
\label{alg:decomposition}
\begin{algorithmic}[1]
\Require Samples $S_1$ with $|S_1| = n$ and samples $S_2$ with $|S_2| = m$, projected onto the SOM space
\Ensure Reduced transport plan $\tilde{P}$ and normalized importance scores $\tilde{I}$ summarizing cluster-level reorganization between $S_1$ and $S_2$
\State Compute the optimal transport plan $P \in \mathbb{R}^{n \times m}$ between $S_1$ and $S_2$ using their BMU locations in the projected SOM space, where $P_{ij}$ denotes the mass transported from sample $i \in S_1$ to sample $j \in S_2$. Use normalized mass if $n\neq m$.
\State Apply HDBSCAN to the projected SOM locations of $S_1$ and $S_2$ separately, yielding clusters $C_a$ ($a \in \{1, \dots, k_1\}$) and $D_b$ ($b \in \{1, \dots, k_2\}$); exclude points labeled as noise
\State Compute centroids $c_a$ for each $C_a$ and $d_b$ for each $D_b$ as the mean SOM prototype of their member nodes
\For{each cluster pair $(a, b)$}
    \State $\tilde{P}_{ab} \gets \sum_{i \in C_a} \sum_{j \in D_b} P_{ij}$
    \State $I_{ab} \gets \tilde{P}_{ab} \cdot \mathrm{MAE}(c_a, d_b)$ \Comment{importance score }
\EndFor
\State $Z \gets \sum_{a, b} I_{ab}$ \Comment{total importance across all cluster pairs}
\For{each cluster pair $(a, b)$}
    \State $\tilde{I}_{ab} \gets I_{ab} / Z$ \Comment{normalized importance score}
\EndFor
\State \Return $\tilde{P}, \tilde{I}$
\end{algorithmic}
\end{algorithm}

\begin{figure}[p]
    \centering
    \begin{subfigure}{0.49\textwidth}
        \centering
        \includegraphics[width=\textwidth]{AGUBeforeDist.png}
        \caption{Before (historical) distribution}
        \label{fig:agu_before_dist}
    \end{subfigure}
    \hfill
    \begin{subfigure}{0.49\textwidth}
        \centering
        \includegraphics[width=\textwidth]{AGUAfterDist.png}
        \caption{After (future) distribution}
        \label{fig:agu_after_dist}
    \end{subfigure}

    \vspace{0.75em}

    \begin{subfigure}{0.8\textwidth}
        \centering
        \includegraphics[width=\textwidth]{TransportPlan.png}
        \caption{Optimal transport plan between clusters}
        \label{fig:agu_transport_plan}
    \end{subfigure}
    \caption{Clustering-based decomposition of an example comparison: historical versus SSP5-8.5 March precipitation in the late-century period (2058--2100), conditioned on Bay-Delta wet years. The before (a) and after (b) sample distributions are each represented as a density over the ClimateSOM space and partitioned into clusters, labeled by their share of the probability mass, with markers at the cluster centroids. The optimal transport plan (c) redistributes mass from the before-clusters (circles) to the after-clusters (diamonds). The \emph{historical}
    distributions shown as the same density rendering and the \emph{future} distributions in the dotted outline. The labeled arrows give the dominant cluster transitions and its importance score $\tilde{I}$.}
    \label{fig:agu_decomposition}
\end{figure}

\begin{figure}[p]
    \centering
    \includegraphics[width=\textwidth]{centroid_weights.pdf}
    \caption{The five highest-scoring cluster transitions (flows) from the transport plan in Figure~\ref{fig:agu_transport_plan}, ranked by importance score. Each column is one transition; the top row shows the source (historical) cluster centroid and the bottom row the target (future) cluster centroid, both rendered as log-standardized precipitation anomaly patterns. The annotations report the redistributed mass, the MAE between the two centroids ($\mathrm{Dist}$), and the resulting normalized importance score. Rendering the centroids directly turns each ranked transition into an interpretable before-and-after pair of spatial precipitation patterns.}
    \label{fig:agu_centroid_weights}
\end{figure}

% \subsection{Conclusion}
\subsection{Conclusion}
\label{subsec:agu-real-conclusion}
This extension carried ClimateSOM from exploration to inference.
%
We equipped the workflow with an EMD-based permutation test that asks, with statistical rigor, whether the distribution of conditioned spatial precipitation patterns has shifted under a changing climate, and we validated it against the standard Welch t-test, finding broad agreement together with an informative asymmetry: the distributional test detects reorganizations of spatial pattern that a mean-based test did not detct.
%
Applied to the co-variability between the Sacramento Bay-Delta Watershed and U.S. West Coast precipitation, the test reveals a statistically significant, regime-dependent decoherence that strengthens with emissions and into the late century, and that is coherently expressed within individual water years rather than emerging only in the period average.
%
The clustering-based decomposition then begins to turn these detected shifts into interpretable spatial pattern transitions, as in the worked example where wet-year March precipitation is projected to redistribute from Southern California toward the north.

% We regard this as a successful sarlitmus test for ClimateSOM.
%
Confronted with a concrete, climatologically substantive question and with the established tools of the domain, the workflow delivers an inferential result those tools were not designed to provide, and renders that result as legible spatial patterns rather than a single summary statistic.
%
The decomposition remains a work in progress, and our broader aim is to fold this inferential and interpretive capability back into the interactive ClimateSOM workflow, so that hypothesis testing and pattern attribution become steps an analyst can take within the same visual environment used to explore the ensemble in the first place.



